\chapter{The elite and non-elite cemeteries}


Following the methodological structure established in the previous chapter, this chapter turns to the non-elite cemeteries of the Napatan period. It employs a dual analytical approach: first, a chronological examination of tomb structure, grave goods, and amulets to trace temporal patterns; and second, a regional analysis to assess the correlation between geographic locations in the Nubian territory and the adoption of certain amulets and use of local and imported materials.	Initial analysis, however, revealed that some sites were not socially homogeneous, but contained a range of tomb types, from simple burial pits to more elaborate structure indicative of elite status. In order to facilitate a precise social and chronological analysis, the elite tombs will be examined in a dedicated section, which will also employ a chronological framework.


\section{The sites}


The sites analyzed in this section are Meroe West and South, Missiminia, Sanam, Mirgissa, Qustul and Debeira 176. The latter three are located in Lower Nubia, between the First and Second Cataracts. The remaining sites are in Upper Nubia: Missiminia lies between the Second and Third Cataracts; Sanam is situated near the royal cemeteries of El-Kurru and Nuri; and Meroe West and South are located further south, beyond the Fifth Cataract. Chronologically, these sites span the entire Napatan period.


Among these sites, only Sanam has direct evidence of an associated settlement \parencites{tuckerDefiningCityNapatan2019}. At Missiminia, potential settlements near the cemetery were likely destroyed by later occupations \parencites[p.~15]{vilaNécropoleMissiminiaSépultures1980}. Sanam represents a particularly significant case, as it was an active trading post contemporary with its cemetery. This is exemplified by a warehouse, known as the 'Treasury', first discovered by Griffith and further excavated in recent years. The excavations revealed raw local and foreign materials, such as elephant tusks and lapis lazuli, alongside finished products including beads and amulets \parencites{griffithOxfordExcavationsNubia1922, vincentelliLongDistanceTrade2018}. This evidence positions Sanam as a central point in long-distance and internal exchange networks. Furthermore, Sanam's population was undoubtedly integrated into this environment and was likely composed of individuals with a variety of professions. The site's proximity with the royal burial sites also suggests that its inhabitants may have been involved with the construction and furnishing of these royal tombs, further suggesting they were skilled workers.


\section{Chronological issues and division}

Dating non-royal tombs is inherently complex, as chronologies primarily rely on inscribed material bearing royal names and comparisons to tombs of fixed date, which can only provide a broad range of possible dates. This makes the precise dating of individual tombs particularly challenging. Consequently, the assigned dates for most of the tombs in the non-royal corpus represent broad time spans that often encompass multiple phases of the Napatan period.


The earliest cemetery, Debeira 176, is considered as contemporaneous with the ancestral El-Kurru burials based on shared features, including tumuli superstructures, contracted body positions, offering niches containing ritually killed vessels, and similar grave goods \parencites[p.~200-202]{säve-söderberghScandinavianJointExpedition1989a}[p.~309-309]{törökTwoWorldsFrontier2009}[p.~631]{lohwasserNapatanNecropoleisBurial2019}. Missiminia also includes burials dated to the pre-25th Dynasty, as Kilroe recently re-examined the excavation report and suggested that four tombs contained post-New Kingdom material \parencites[p.~251]{kilroePilgrimFlasksChronological2019}. Furthermore, the earlier tombs at Missiminia display similarities with the post-New Kingdom Cemetery D in Amara West \parencites[p.~57]{binderCemeteryAmaraWest2011}. However, the majority of the tombs at Missiminia are dated from the 25th Dynasty to the Early Napatan period, due to similarities with Sanam and Qustul \parencites{williamsExcavationsAbuSimbel1990} [p.~300-301]{lohwasserAspekteNapatanischenGesellschaft2012}. The chronology of the other sites varies: Mirgissa is dated to the 25th Dynasty, while Meroe West and South cemeteries have the longest use, with phases spanning from the 25th Dynasty into the Late Napatan period \parencites{dunhamWestSouthCemeteries1963, geusCimetièreXXVèmeDynastie1975}. Sanam, however, has enough evidence to conclude that it was founded during the New Kingdom occupation, when cave and stair tombs were built, but the main use of the cemetery was during the 25th Dynasty-Early Napatan phases \parencites{lohwasserAspekteNapatanischenGesellschaft2012}.


\section{Elite tombs}


\subsection{Tomb structure}


Architectural variations at the cemeteries of Meroe West, South and Sanam points to heterogeneous social groups. This is evidenced by super- and substructures---such as pyramids, chambers, stair and cave graves---, which bear similarities to royal tombs, indicating that elite individuals were also buried there (Figure \ref{fig:elite-tombs}). In total, there are 280 tombs belonging to the elite, dating from the 25th Dynasty to the Late Napatan period. Only two elite tombs in Meroe West were dated by Reisner/Dunham to the 25th Dynasty, therefore, for a better statistical analysis, they were included in the 25th Dynasty-Early Napatan phase.


Analysis of the tombs reveals that pyramids are only present at Meroe West and South, with pits and single chambers as the substructures. Sanam, on the other hand, contained cave tombs built in the compact soil, with chambers (up to four), but no identifiable superstructure. Access to these tombs was possible by either stairs, pits or slopes. All tombs are dated to the overall span of the Sanam cemetery, from the 25th Dynasty to the Early Napatan phase. Overall, pyramids were a consistent feature across all phases, although the substructure evolved from pits to chambers in later tombs. Additionally, graves with unidentified superstructures, were exclusively paired with chambers and cave tombs.


\begin{figure}[H]
	\centering
	\includegraphics[width=1\textwidth]{elite-tombs}
	\caption{Distribution of super- and substructures in elite tombs by chronological phase.}
	\label{fig:elite-tombs}
\end{figure}


Architectural features reveal further similarities between elite and royal burials (Figure \ref{fig:elite-tomb-elements}). Holes for burial beds, for instance, a feature found only in earlier tombs, were replaced by coffin benches in later burials. Decoration was rare and exclusive to a few Meroe South tombs from the Late Napatan phase. Pyramid S 8, for example, featured decorated chapel blocks with Egyptian hieroglyphics, while S 7 contained chapel wall reliefs. Chapels are present in all phases in the Meroe cemeteries, most likely due to better preservation conditions, but are more numerous in early and later tombs, with Middle Napatan burials containing the fewest. While stairs appear in every phase, they are least common in Early Napatan tombs. Furthermore, the higher status of the Meroe cemeteries, in particular, is underscored by the presence of links to royalty. Three significant tombs---W 609, S 85, and S 155---are associated with royal wives, who, for an unknown reason, were not buried at El-Kurru and Nuri \parencites[p.~107]{lohwasserKöniglichenFrauenIm2001}. Notably, pyramid S 85 yielded inscriptions with the name Mernua. The chronology of these elite burials spans from the 25th Dynasty (W 609 and S 155) to the Early Napatan phase (Mernua's tomb). The tomb of Mernua is very similar to contemporary royal tombs, especially of wives, containing a stairway, a single chamber and a coffin bench. However, it did not include a chapel, nor any sign of decoration, which could be due to preservation conditions.


\begin{figure}[H]
	\centering
	\includegraphics[width=1\textwidth]{elite-tomb-elements}
	\caption{Distribution of tomb features in elite burials by chronological phase.}
	\label{fig:elite-tomb-elements}
\end{figure}


\subsection{Object types}


\begin{figure}[H]
	\centering
	\includegraphics[width=1\textwidth]{elite-objs}
	\caption{Distribution of object types in elite tombs by chronological phase.}
	\label{fig:elite-objs}
\end{figure}

Overall, object types at elite tombs are rather similar to royal assemblages, with beads, vessels and adornments as the most common categories across all phases (Figure \ref{fig:elite-objs}). However, one difference is that shabtis were much rarer in elite contexts. Most of them were found in Mernua's tomb (Early Napatan), reinforcing her connection with the royal sphere. The others were mostly fragments coming from both Early and Middle Napatan tombs. One, from tomb S 82, included an inscription with a name similar to king's Malaqen, but with a different spelling and a female determinative. On the other hand, the fragments from tomb S 84 belonged to three different people. Other shabtis were found in different tombs, including non-elite tombs and the elite burials S 42 and S 49. For that reason, Dunham points out that the shabtis were so scattered by the looting that it was not possible to determine their tomb of origin \parencites[p.~406-409, 4015-416]{dunhamWestSouthCemeteries1963}. Consequently, it seems that Mernua's intact burial was the only one which contained a complete set of 88 shabtis. Her tomb also yielded all of the heart scarabs and canopics.


Games and toys are more present in elite contexts than royal ones, where they were only found in El-Kurru during the 25th Dynasty, although this could be due to preservation issues if these objects were made of perishable materials. In contrast, at Meroe South, they were found in burials from the 25th Dynasty to Middle Napatan, with the highest concentration in tomb S 86 that is dated to either the 25th Dynasty or Early Napatan phases. Similarly, staffs---previously only found in Nuri royal tombs from the 25th Dynasty to the Late Napatan period---, are only found in Late Napatan elite burials at Meroe South.


Burial containers were present in most phases. The data reveals a general pattern: coffins were the predominant type, and their distribution varied chronologically and geographically. The earliest evidence comes from a single coffin in an 25th Dynasty-Early Napatan burial at Meroe West. This is followed by one coffin from Early Napatan (Meroe South), two from Early-Middle Napatan (Meroe West), and four from Middle Napatan (Meroe South). The sole Late Napatan coffin was also found in Meroe South, suggesting either the location was better suited for preservation or people there had more access to burial containers. A significant concentration was found at Sanam, which included three coffins and eight tombs with traces of cartonnage. The unique bed burial was found in a tomb at Meroe South dated to the 25th Dynasty-Early Napatan phases.


Early elite assemblages are not as diverse as royal ones, comprising of only five object types and, particularly, lacking some burial equipments such as canopics, mummy trappings, shabtis, stelae and heart scarabs. One possibility is that, as almost all of these tombs were disturbed, these objects were looted. Cylinder sheaths were absent from elite burials, supporting the idea that they were exclusive to the royal sphere. Moreover, objects already rare in royal tombs, such as shrines, statuettes, pebbles and offering vessels, were also not found with the elite. Headrests are the only object type exclusive to elite assemblages, although only one ivory example was found in a cave tomb from Sanam.


\subsection{Amulets}


\subsubsection{Amulet distribution}


Amulets from the elite sphere are rarer than those from royalty, although looting has to be considered. Out of the total 181 tombs, only 29 contained amulets (Table \ref{tab:elite-distro}). In absolute numbers, tombs from the 25th Dynasty to Early Napatan contained more amulets, with Mernua's tomb being the only one from the Early Napatan phase to contain amulets, totaling 64.
The highest percentage of tombs with amulets was from the Late Napatan phase (17\%), while the previous phases barely reached 15\%. Overall, very few elite tombs contained amulets; the median for all phases was 0, as the majority of tombs did not contain any amulets. Furthermore, amulets from elite contexts follow the same pattern as those from the royal sphere, with fewer amounts in later phases.


{\small
	\begin{xltabular}{\textwidth}{X X X X X}
		\caption{Distribution of amulets in elite tombs by phase, including total counts of amulets and tombs, tombs containing amulets and amulet frequency (\%).}
		\label{tab:elite-distro} \\
		\toprule
		\textbf{Chronological phase} & \textbf{Total number of tombs} & \textbf{Total number of amulets} & \textbf{N. of tombs with amulets} & \textbf{\% of tombs with amulets} \\
		\midrule
		\endfirsthead
		\multicolumn{5}{c}{\textit{Continued from previous page}} \\
		\toprule
		\textbf{Chronological phase} & \textbf{Total number of tombs} & \textbf{Total number of amulets} & \textbf{N. of tombs with amulets} & \textbf{\% of tombs with amulets} \\
		\midrule
		\endhead
		\bottomrule
		\endlastfoot
		25th Dynasty--Early Napatan & 174 & 86 & 18 & 10\% \\
		\cmidrule(r){1-5}
		Early Napatan & 9 & 64 & 1 & 11\% \\
		\cmidrule(r){1-5}
		Early Napatan--Middle Napatan & 7 & 14 & 1 & 14\% \\
		\cmidrule(r){1-5}
		Middle Napatan & 43 & 9 & 5 & 12\% \\
		\cmidrule(r){1-5}
		Middle Napatan--Late Napatan & 21 & 1 & 1 & 5\% \\
		\cmidrule(r){1-5}
		Late Napatan & 18 & 7 & 3 & 17\% \\
		\midrule
		\textbf{Total} & \textbf{281} & \textbf{181} & \textbf{29} & -- \\
	\end{xltabular}
}


\subsubsection{Amulet types}

\begin{figure}[H]
	\centering
	\includegraphics[width=1\textwidth]{elite-types}
	\caption{Distribution of amulet types in elite tombs by chronological phase.}
	\label{fig:elite-types}
\end{figure}

Continuing the trend of similarities with royal contexts, amulets of symbols are the most common category in elite tombs, with the exception of the single Middle-Late Napatan tomb, which contained a single animal amulet (Figure \ref{fig:elite-types}). Deities are the second most popular amulet. however, they were entirely absent from tombs dated to the Middle-Late to Late Napatan range, which are also the phases with less amulets. Animals are the only amulet category present in all phases, being more prevalent in earlier tombs, similarly to deities and symbols. Amulets depicting humans is a type that is first attested in elite contexts, although only one example was found in a cave tomb from Sanam. Amulets depicting nature elements and objects were also rare in elite burials.


\subsubsection{Amulet motifs}

\begin{figure}[H]
	\centering
	\begin{subfigure}{1\textwidth}
		\centering
		\includegraphics[width=\linewidth, keepaspectratio]{elite-motifs-1}
	\end{subfigure}
	
	\begin{subfigure}{1\textwidth}
		\centering
		\includegraphics[width=\linewidth, keepaspectratio]{elite-motifs-2}
	\end{subfigure}
	\caption{Distribution of amulet motifs in elite tombs by chronological phase.}
	\label{fig:elite-motifs}
\end{figure}


The \textit{udjat} is the most widespread amulet motif, especially in earlier phases ({Figure \ref{fig:elite-motifs}}). Mernua's tomb from the Early Napatan phase yielded the greatest variety of motifs. Most of the motifs are unique and were cut from sheet electrum (Figures \ref{fig:elite-collar}, \ref{fig:elite-winged}, \ref{fig:elite-mernua-pics}). Among them are those representing hieroglyphic signs, such as the \textit{sema} (“unification”), \textit{sa} (“protection”) and \textit{akhet} (“horizon”), and as well as symbolic motifs, some related to royalty, namely the \textit{ankh}, whip, \textit{was} scepter and sun boat. The tomb also contained the only amulet in the corpus of the god Aker, the personification of the horizon, depicted as two recumbent lion torsos merged with the sun disc; as well as the only \textit{menat}, represented differently than the examples from the royal burials, without further decoration (Figure \ref{fig:elite-mernua-pics}). Additionally, the sheet amulets are similar in style to gold royal examples from the 25th Dynasty to the Middle Napatan period.


Especially noteworthy are two amulets depicting winged animals (Figure \ref{fig:elite-winged}). One depicts a winged snake, the sole piece in the corpus. The other amulet is of a winged griffin, a motif also found in the tomb of the Late Napatan king Nastasen (Nu. 15). Although both amulets are not stylistically equal, both are made of sheet metal and depict a winged quadruple animal with an avian head, which could be interpreted as a griffin, as it does not resemble any other animals. The major differences between both amulets are their distant dating and location, indicating the rarity of this motif in Napatan Nubia. As mentioned in the previous chapter, in Egypt, griffins have been interpreted in different ways depending on the context. For Napatan Nubia, as there are only these two amulets in this fashion in very different contexts, it is not clear which symbolic meaning was attached to the griffin. Nevertheless, Mernua's amulets are further evidence of her link to the royal sphere, although probably in a position lesser than minor wives, as she was not interred in the royal burial sites.

\begin{figure}
	\centering
	\includegraphics[width=1\textwidth, height=4cm, keepaspectratio]{elite-23.1435a-d}
	\caption{From Mernua's tomb, amulets of a vulture, \textit{uraeus} and sun boat (MFA 23.1435a-d).}
	\label{fig:elite-collar}
\end{figure}

\begin{figure}[H]
	\centering
	\begin{minipage}{0.49\textwidth}
		\centering
		\includegraphics[width=\textwidth, height=3cm, keepaspectratio]{elite-23.1430}
	\end{minipage}
	\hspace{-53mm}%
	\begin{minipage}{0.49\textwidth}
		\centering
		\includegraphics[width=\textwidth, height=3cm, keepaspectratio]{elite-23.1431}
	\end{minipage}
	\caption{Amulets of a winged snake and griffin from the tomb of Mernua (MFA 23.1430, 23.1431).}
	\label{fig:elite-winged}
\end{figure}


A few motifs were only found in amulets from Middle Napatan burials onwards---hawk, falcon and the goddess Sekhmet. For Early to Middle Napatan tombs, exclusive motifs include flies, headrests and two unidentified deities, one with a snake head and another with a hawk head. There was a single fly amulet, as well as a headrest, both from the Early-Middle Napatan range. Moreover, amulets of Thoth were only found from the Early to Middle Napatan phases. Interestingly, amulets of this god were not found at the royal cemeteries. Along with the the snake and hawk-headed unidentified deities and Sekhmet, these were the only amulet depicting deities present in Early-Middle Napatan and later burials. Egyptian-style amulets continue to be popular among elite burials, specially in Mernua's tomb. However, some amulets appear to be exclusive to royal contexts, such as the \textit{menat}, the Four Sons of Horus and hieroglyphic signs---motifs only present with Mernua.

\begin{figure}[H]
	\centering
	\begin{subfigure}{0.25\textwidth}
		\includegraphics[width=\linewidth, height=3cm, keepaspectratio]{elite-23.1425}
	\end{subfigure}%
	\hspace{-5mm}%
	\begin{subfigure}{0.25\textwidth}
		\includegraphics[width=\linewidth, height=3cm, keepaspectratio]{elite-23.1409}
	\end{subfigure}%
	\hspace{-20mm}%
	\begin{subfigure}{0.25\textwidth}
		\includegraphics[width=\linewidth, height=3cm, keepaspectratio]{elite-23.1406}
	\end{subfigure}%
	\\[0mm]
	\begin{subfigure}{0.25\textwidth}
		\includegraphics[width=\linewidth, height=3cm, keepaspectratio]{elite-23.1414}
	\end{subfigure}%
	\hspace{-17mm}%
	\begin{subfigure}{0.25\textwidth}
		\includegraphics[width=\linewidth, height=3cm, keepaspectratio]{elite-23.1446}
	\end{subfigure}%
	\hspace{-15mm}%
	\begin{subfigure}{0.25\textwidth}
		\includegraphics[width=\linewidth, height=3cm, keepaspectratio]{elite-23.1419}
	\end{subfigure}%
	\hspace{-23mm}%
	\begin{subfigure}{0.25\textwidth}
		\includegraphics[width=\linewidth, height=3cm, keepaspectratio]{elite-23.1420}
	\end{subfigure}%
	\hspace{-23mm}%
	\begin{subfigure}{0.25\textwidth}
		\includegraphics[width=\linewidth, height=3cm, keepaspectratio]{elite-23.1418}
	\end{subfigure}
	\caption{Amulets from Mernua's tomb: the god Aker (MFA 23.1425), \textit{menat} (MFA 23.1409), \textit{sema} (MFA 23.1406), \textit{ankh} (MFA 23.1414), \textit{akhet} (MFA 23.1446), a whip (MFA 23.1419), a \textit{was} scepter (MFA 23.1420) and \textit{sa} sign (MFA 23.1418)}
	\label{fig:elite-mernua-pics}
\end{figure}

\subsubsection{Amulet materials}

\begin{figure}[H]
	\centering
	\includegraphics[width=1\textwidth]{elite-materials}
	\caption{Distribution of amulet materials in elite tombs by chronological phase.}
	\label{fig:elite-materials}
\end{figure}

The material composition of elite amulets differ from royal ones in a particular way: faience was the most common material only during the 25th Dynasty-Early Napatan phase (Figure \ref{fig:elite-materials}). Although appearing in burials spanning from the Early to Middle Napatan phases, they were absent in later burials and from the only Early Napatan tomb with amulets (Mernua's). This tomb included amulets made from a variety of materials, mostly electrum, but not a single one from faience, which makes it exceptional, as even the royal burials contained many faience amulets. 

Amulets made of gold were also rare, as only five were found, which could be a result of robbery. Materials sourced from the Egyptian Eastern Desert---such as jasper, serpentinite, travertine, carnelian and chalcedony---, were used throughout all phases, albeit in small numbers. Locally sourced materials, namely haematite and agate, were rarer. Not surprising, Mernua's amulets were the ones composed of the most foreign materials, which can be due to her social status.


Sandstone and steatite were the only elite materials not found in royal amulets, although this material was used for other grave goods, such as shabtis and scarabs. This evidence may indicate they were not popular materials for royal amulets, specifically. Regardless, only two 25th Dynasty-Early Napatan elite amulets were made of steatite. Although, it has no known ancient quarries, some likely existed in the Eastern Desert \parencites[p.~384]{harrellArchaeologyGeologyAncient2024}. Another pair of amulets from the same Early-Middle Napatan tomb, were made with sandstone, which was sourced from the Nile Valley and near Aswan \parencites[p.~506]{harrellArchaeologyGeologyAncient2024}. The only material not sourced within the Egyptian-Sudanese territory is lapis lazuli, which accounts for nine amulets, eight in a single 25th Dynasty tomb (W 486) and one in a 25th Dynasty-Early Napatan burial from Sanam, reinforcing their elite status. Although, from the Middle Napatan period onwards the number of amulets decreases, materials coming from the Egyptian Eastern Desert were still being used in their manufacture.


\section{Non-elite tombs}


\subsection{Pre-25th Dynasty}


\subsubsection{Tomb structure}

\begin{figure}[H]
	\centering
	\includegraphics[width=1\textwidth]{pre-tombs}
	\caption{Distribution of super- and substructures in pre-25th Dynasty non-elite tombs.}
	\label{fig:pre-tombs}
\end{figure}


Pre-25th Dynasty non-elite tombs are located at two sites: Debeira 176 and Missiminia (Figure \ref{fig:pre-tombs}). At Debeira 176, a site consisting entirely of tombs from this phase, tumuli are the sole superstructure, similar to contemporary royal tombs from El-Kurru. In contrast, the superstructures of the four tombs at Missiminia are unidentified, being either absent or eroded \parencites[p.~170]{vilaNécropoleMissiminiaSépultures1980}. Substructures in both sites consisted entirely of burial pits.


All four Missiminia tombs contained side niches. Half were lined with bricks, including roofs, while two included trenches---likely for a bed burial, although none were found. At Debeira 176, tomb features were less consistent. Of the total 104 tombs, 35 were roofed, although not always with bricks, and 20 included niches or offering places---a feature that also connects these tombs with the ancestral royal burials. Trenches were found in a single tomb.


\subsubsection{Object types}


\begin{figure}[H]
	\centering
	\includegraphics[width=1\textwidth, height=8cm, keepaspectratio]{pre-objs}
	\caption{Distribution of object types in pre-25th Dynasty non-elite tombs.}
	\label{fig:pre-objs}
\end{figure}


Object distribution in pre-25th Dynasty non-elite tombs is very similar to contemporaneous royal ones, with beads, adornments and ceramic vessels being the most common (Figure \ref{fig:pre-objs}). However, non-elite assemblages are less diverse, lacking such items as weapons, jar stoppers, shabtis, furniture and offering tables and vessels. As Missiminia had only four tombs, it is not surprising that they featured only three object types. Overall, the non-elite assemblage contains objects that were not necessarily made for the burial, but could have been used in daily life.


\subsubsection{Amulets}


\paragraph{Amulet distribution}


There were only eleven amulets from this phase, distributed across five tombs (Table \ref{tab:pre-distro}). Nevertheless, the distribution bares similarities to contemporary royal contexts, which also included a low number of amulets. Out of the 104 tombs from Debeira 176, only 4 contained at least one amulet, 5\%. At Missinimia, this percentage is higher (25\%). However, this is due to only one tomb out of four containing a single amulet.


{\small
	\begin{xltabular}{\textwidth}{X X X X X}
		\caption{Distribution of amulets in pre-25th Dynasty non-elite tombs by site, including total counts of amulets and tombs, tombs containing amulets and amulet frequency (\%).}
		\label{tab:pre-distro} \\
		\toprule
		\textbf{Site} & \textbf{Total number of tombs} & \textbf{Total number of amulets} & \textbf{N. of tombs with amulets} & \textbf{\% of tombs with amulets} \\
		\midrule
		\endfirsthead
		\multicolumn{5}{c}{\textit{Continued from previous page}} \\
		\toprule
		\textbf{Site} & \textbf{Total number of tombs} & \textbf{Total number of amulets} & \textbf{N. of tombs with amulets} & \textbf{\% of tombs with amulets} \\
		\midrule
		\endhead
		\bottomrule
		\endlastfoot
		Debeira 176 & 104 & 10 & 4 & 5\% \\
		\cmidrule(r){1-5}
		Missiminia & 4 & 1 & 1 & 25\% \\
		\midrule
		\textbf{Total} & \textbf{108} & \textbf{11} & \textbf{5} & -- \\
	\end{xltabular}
}


\paragraph{Amulet types}


In contrast with the contemporary royal contexts, where most amulets were of symbols, deities are the most common type in non-elite graves, with only one instance of a symbol (Figure \ref{fig:pre-types}). In fact, animals are the second most common, with three instances from Debeira 176 and the sole amulet from Missiminia. Other types of amulets were absent from these tombs.


\begin{figure}[H]
	\centering
	\includegraphics[width=1\textwidth]{pre-types}
	\caption{Distribution of amulet types in pre-25th Dynasty non-elite tombs.}
	\label{fig:pre-types}
\end{figure}


\paragraph{Amulet motifs}

\begin{figure}[H]
	\centering
	\includegraphics[width=1\textwidth]{pre-motifs}
	\caption{Distribution of amulet motifs in pre-25th Dynasty non-elite tombs.}
	\label{fig:pre-motifs}
\end{figure}


Regarding motifs, unfortunately, the majority of the deity amulets from Debeira 176 are mostly unidentified, as they were very damaged (Figure \ref{fig:pre-motifs}). But the site did include an amulet of the god Bes, as well as an \textit{udjat}, and three lions. The sole amulet from Missiminia depicts a ram's head, similar to those found in royal tombs from Middle Napatan. However, this particular one featured a single \textit{uraeus}, instead of two (Figure \ref{fig:pre-ram-head}). The inclusion of lion amulets is interesting, since this motif was absent in royal contexts and only one was found in an elite burial from the 25th Dynasty-Early Napatan phase.


\begin{figure}[H]
	\centering
	\includegraphics[width=1\textwidth, height=4cm, keepaspectratio]{pre-ram-head}
	\caption{Ram's head amulet \parencites[p.~112, fig.~118, 2]{vilaNécropoleMissiminiaSépultures1980}}
	\label{fig:pre-ram-head}
\end{figure}


\paragraph{Amulet materials}


The eleven pre-25th Dynasty amulets were made of only three materials (Figure \ref{fig:pre-materials}). Faience and carnelian, equally represented, comprised 90\% of the corpus and were found exclusively at Debeira 176. In contrast, the single amulet from Missiminia was made of gold, mirroring the contemporary royal preference for this material. Furthermore, the presence of carnelian is interesting, as this material was likely sourced from quarries near Aswan, suggesting access to this particular material.


\begin{figure}[H]
	\centering
	\includegraphics[width=1\textwidth]{pre-materials}
	\caption{Distribution of amulet materials in pre-25th Dynasty non-elite tombs.}
	\label{fig:pre-materials}
\end{figure}


\subsection{25th Dynasty}


\subsubsection{Tomb structure}


\begin{figure}[H]
	\centering
	\includegraphics[width=1\textwidth]{25-tombs}
	\caption{Distribution of super- and substructures in 25th Dynasty non-elite tombs.}
	\label{fig:25-tombs}
\end{figure}


While tumulus was the most common superstructure in the previous phase, in the 25th Dynasty, most of the 154 tombs had no superstructure (Figure \ref{fig:25-tombs}). All tombs from Meroe West and South were either unidentified or had no superstructure at all. Mirgissa presents an interesting case, as the majority of superstructures there were noted as having disappeared, which could mean they existed at one point. As for substructures, pits continued to be the only type.


Differences between the tombs at the Meroe region and Mirgissa also show in other architectural features. Notably, burial bed cutouts were found in 34 graves at Meroe but none at Mirgissa, where all 27 pits were lined with bricks. Uncommon features include roofing, present in one tomb per site (W 611, S 216 and Mirgissa 26), and trenches in five Meroe tombs, potentially for grave goods or burial bed supports; although, no traces of burial beds were found.


\subsubsection{Object types}


25th Dynasty non-elite object assemblages vary significantly by location. While in the Meroe region there is a greater diversity of grave goods, in Mirgissa, grave goods included only adornments, beads and scarabs (Figure \ref{fig:25-objs}). 28\% of the total tombs did not contain any objects, the majority at Mirgissa. Beads, ceramic and non-ceramic vessels are among the most common object types at the Meroe sites. Scarabs, although present in all three sites, were found in greater numbers at Meroe West---with tomb W 609 accounting for 201 instances.


Two categories of objects are exclusive to Meroe West graves: a single saucer lamp and five seals. While all seals contained engraved symbols, one from W 658 was inscribed with cartouches of Kashta and Amenirdis on one side, while the other included an \textit{udjat} \parencites[p.~304]{dunhamWestSouthCemeteries1963}. Another from W 671 included Tutmosis III's cartouche (MFA 23.858). Pebbles, previously only found in two 25th Dynasty royal burials, were found in four non-elite tombs, totaling five instances, with the majority at Meroe West. Fragments of shabtis were found in five tombs at Meroe South. The ones from S 88, S 125 and S 214 bore a woman's name and were also found in a tomb from the Early Napatan phase. As noted in the report, the shabtis were likely scattered by the looters and their tomb of origin could not be identified \parencites[p.~416]{dunhamWestSouthCemeteries1963}. These disperse fragments account for 33 of the total 35 shabtis found in 25th Dynasty tombs. The remaining two, from S 126 and S 134, were also fragments and, although, not noted in the reports whether they were intrusive, both tombs had been plundered \parencites[p.~357, 429]{dunhamWestSouthCemeteries1963}. Therefore, it is not possible to determine positively that non-elite tombs from this phase contained shabtis.


As for burial containers, two tombs from Meroe West and South contained eye inlays which could have been part of a wooden coffin that was not preserved. However, none of the tombs contained a coffin bench or cutouts for a burial bed. Bead nets were also more common in non-elite contexts than royal and elite ones. A tomb from Meroe West contained the only offering table from this phase. Games and toys were also present in this contexts, with 19 instances from the Meroe cemeteries.


\begin{figure}[H]
	\centering
	\includegraphics[width=1\textwidth, height=11cm, keepaspectratio]{25-objs}
	\caption{Distribution of object types in 25th Dynasty non-elite tombs.}
	\label{fig:25-objs}
\end{figure}


\subsubsection{Amulets}


\paragraph{Amulet distribution}


A total of 1577 amulets were distributed across 62 tombs (Table \ref{tab:25-distro}). At Meroe West, which contained 94\% of the amulets, over half of the tombs (54\%) contained at least one amulet. In contrast, at Meroe South, 27\% of tombs included amulets, while Mirgissas's percentage was even lower at 22\%.


{\small
	\begin{xltabular}{\textwidth}{X X X X X}
		\caption{Distribution of amulets in 25th Dynasty non-elite tombs by site, including total counts of amulets and tombs, tombs containing amulets and amulet frequency (\%).}
		\label{tab:25-distro} \\
		\toprule
		\textbf{Site} & \textbf{Total number of tombs} & \textbf{Total number of amulets} & \textbf{N. of tombs with amulets} & \textbf{\% of tombs with amulets} \\
		\midrule
		\endfirsthead
		\multicolumn{5}{c}{\textit{Continued from previous page}} \\
		\toprule
		\textbf{Site} & \textbf{Total number of tombs} & \textbf{Total number of amulets} & \textbf{N. of tombs with amulets} & \textbf{\% of tombs with amulets} \\
		\midrule
		\endhead
		\bottomrule
		\endlastfoot
		Meroe South & 48 & 37 & 13 & 27\% \\
		\cmidrule(r){1-5}
		Meroe West & 79 & 1490 & 43 & 54\% \\
		\cmidrule(r){1-5}
		Mirgissa & 27 & 50 & 6 & 22\% \\
		\midrule
		\textbf{Total} & \textbf{154} & \textbf{1577} & \textbf{62} & -- \\
	\end{xltabular}
}


\paragraph{Amulet types}


\begin{figure}[H]
	\centering
	\includegraphics[width=1\textwidth]{25-types}
	\caption{Distribution of amulet types in 25th Dynasty non-elite tombs.}
	\label{fig:25-types}
\end{figure}


The amulet assemblages from Meroe West and Mirgissa are the most diverse, featuring all of the seven different types (Figure \ref{fig:25-types}). The distribution of these types, however, varies between the sites. At Meroe West, the distribution follows a general pattern, with symbols being the most prevalent category, followed by deities and animals. Mirgissa presents a different arrangement, with deities as the dominant type, followed by symbols and then animals.


Meroe South exhibits a more balanced distribution. While symbols and deities are the most common there, the difference in numbers is less pronounced. Notably, amulets of humans and other motifs were slightly more common than animal amulets, and objects and nature elements were entirely absent. Furthermore, the 25th Dynasty non-elite tombs contain a significant quantity of a new amulet type previously found in a single elite 25th Dynasty-Early Napatan burial: human figures and body parts, present in all three sites.


\paragraph{Amulet motifs}


\begin{figure}[H]
	\centering
	\includegraphics[width=1\textwidth, height=12cm, keepaspectratio]{25-deity-motifs}
	\caption{Distribution of deity motifs in 25th Dynasty non-elite tombs.}
	\label{fig:25-deity-motifs}
\end{figure}


Amulets of deities depict 18 different ones, with the most variety from Meroe West (Figure \ref{fig:25-deity-motifs}). The Pataikos is the most common in that site, while Sekhmet was the most popular in Mirgissa and Meroe South. Eight amulets were of unknown deities due to their fragmented nature, making a definite identification impossible. The Pataikos amulets displayed different designs, with one featuring two round gold adornments under his armpits (Figure \ref{fig:25-pataikos}). Meroe West was home to the first amulets in the corpus of the goddesses Taweret and Nephthys, and the gods Horus (not in his child form) and Nefertum. Other popular motifs include the goddesses Mut, Hathor, Bastet, Isis nursing Horus, and the god Bes. In contrast, other deities were less popular, such as Isis, Shu, Bastet, Horus child, Anubis, Amun, and Ra-Horakhty.


One amulet depicting a group of three deities from Meroe South, unfortunately, did not include their identification and the reliefs are unclear (Figure \ref{fig:25-three-deities}). However, the deity on the left is possibly Isis or Hathor due to the crown featuring cow horns and full moon. An amulet depicting a lion-headed goddess suckling a child is possibly of Bastet and Horus, as the goddess was represented as nursing a king in the enthronement ceremony in Harsyotef and Nastasen stelae (Figure \ref{fig:25-lion-nursing}) \parencites[p.~182]{kormyshevaGodsDivineSymbols2006}. If this interpretation is correct, this would mean the presence of royal iconography in a non-elite context. An amulet from S 125 also features a seated lion-headed goddess with hieroglyphic signs on the back (Figure \ref{fig:25-openwork}). The design is particular interesting as all sides, except for the back, are openwork, possibly serving as the suspension method.


 \begin{figure}[H]
	\centering
	\begin{minipage}{0.30\textwidth}
		\centering
		\includegraphics[width=\textwidth, height=4cm, keepaspectratio]{25-23.827}
	\end{minipage}
	\hspace{-25mm}%
	\begin{minipage}{0.30\textwidth}
		\centering
		\includegraphics[width=\textwidth, height=4cm, keepaspectratio]{25-23.844}
	\end{minipage}
	\hspace{-18mm}%
	\begin{minipage}{0.30\textwidth}
		\centering
		\includegraphics[width=\textwidth, height=4cm, keepaspectratio]{25-47.3887}
	\end{minipage}
	\caption{Different designs of Pataikos amulets (MFA 23.827, 23.844, 47.3887).}
	\label{fig:25-pataikos}
\end{figure}


\begin{figure}[H]
	\centering
	\begin{minipage}{0.30\textwidth}
		\centering
		\includegraphics[width=\linewidth, height=4cm, keepaspectratio]{25-24.996}
		\captionof{figure}{Amulet depicting a group of three deities (MFA 24.996).}
		\label{fig:25-three-deities}
	\end{minipage}%
	\hspace{7pt}
	\begin{minipage}{0.30\textwidth}
		\centering
		\includegraphics[width=\linewidth, height=4cm, keepaspectratio]{25-2179}
		\captionof{figure}{Amulet of a lion-headed goddess nursing a child (SNM 2179).}
		\label{fig:25-lion-nursing}
	\end{minipage}%
	\hspace{7pt}
	\begin{minipage}{0.30\textwidth}
		\centering
		\includegraphics[width=\linewidth, height=4cm, keepaspectratio]{25-24.997}
		\captionof{figure}{Openwork amulet of a seated lion-headed deity (MFA 24.997).}
		\label{fig:25-openwork}
	\end{minipage}
\end{figure}


\begin{figure}[H]
	\centering
	\includegraphics[width=1\textwidth, height=12cm, keepaspectratio]{25-udjat-motifs}
	\caption{Distribution of \textit{udjat} motifs in 25th Dynasty non-elite tombs.}
	\label{fig:25-udjat-motifs}
\end{figure}


25th Dynasty non-elite tombs contained the highest amount of \textit{udjat} amulets, with 874 examples, especially at Meroe West, which accounts for 97\% of the total (Figure \ref{fig:25-udjat-motifs}). These amulets were distributed across 32 tombs, with W 609 alone containing 480 examples, over half of the total. Conversely, both Mirgissa and Meroe South had only 12 \textit{udjat} amulets each. In fact, all symbol amulets from Meroe South were of the \textit{udjat}. 


\begin{figure}[H]
	\centering
	\includegraphics[width=1\textwidth]{25-symbol-motifs}
	\caption{Distribution of symbol motifs (excluding \textit{udjat}) in 25th Dynasty non-elite tombs.}
	\label{fig:25-symbol-motifs}
\end{figure}


Regarding the remaining symbol amulets, the \textit{uraeus} was particularly popular in Meroe West, but only one was found in Mirgissa. Meroe West also included two unique amulets depicting a lion-headed \textit{uraeus} and a double \textit{ba} bird (Figure \ref{fig:25-symbol-motifs}). Apart from the \textit{uraeus}, Mirgissa also included the first example of a quadruple \textit{udjat} amulet. The less diverse assemblage of symbol motifs is a stark contrast from contemporary royal and elite contexts.


\begin{figure}[H]
	\centering
	\includegraphics[width=0.49\textwidth, height=3.5cm, keepaspectratio]{25-2199}
	\hspace{-2mm}%
	\includegraphics[width=0.49\textwidth, height=3.5cm, keepaspectratio]{25-23.822}
	\caption{Amulets of a lion-headed \textit{uraeus} and a double ba bird (SNM 2199, MFA 23.822).}
	\label{fig:25-symbol-pics}
\end{figure}


\begin{figure}[H]
	\centering
	\includegraphics[width=1\textwidth]{25-animal-motifs}
	\caption{Distribution of amulet animal motifs in 25th Dynasty non-elite tombs.}
	\label{fig:25-animal-motifs}
\end{figure}


\begin{figure}[H]
 	\centering
 	\begin{minipage}{0.22\textwidth}
 		\includegraphics[width=\linewidth, height=5cm, keepaspectratio, clip, trim={0cm 1.7cm 0cm 1.7cm}]{25-23.755}
 	\end{minipage}%
 	\hspace{-2mm}%
 	\begin{minipage}{0.22\textwidth}
 		\includegraphics[width=\linewidth, height=5cm, keepaspectratio]{25-23.770}
 	\end{minipage}%
 	\hspace{-2mm}%
 	\begin{minipage}{0.22\textwidth}
 		\includegraphics[width=\linewidth, height=5cm, keepaspectratio]{25-47.3027}
 	\end{minipage}%
 	\hspace{-2mm}%
 	\begin{minipage}{0.23\textwidth}
 		\includegraphics[width=\linewidth, height=5cm, keepaspectratio]{25-23.1459}
 	\end{minipage}
 	\caption{Amulets of two frogs, an ibis, a double headed ram and a lion-headed fly (MFA 23.755, 23.770, 47.3027, 23.1459).}
 	\label{fig:25-animal-pics}
\end{figure}


Motifs of animals were more diverse than those of deities and symbols, with 20 different designs, 14 of those were exclusive to Meroe West (Figure \ref{fig:25-animal-motifs}). Furthermore, the corpus is also more diverse than both royal and elite. The most common motif were 26 lion-headed flies, exclusive to non-elite contexts, 25 from tomb W 608 and one from W 861 (pictured last in Figure \ref{fig:25-animal-pics}). A single tomb, W 603, featured the 16 total standard fly amulets. This is an increase in the number of fly amulets compared to royal and elite assemblages.


The second most popular amulets were rams, present in all three sites. Additionally, designs featuring a double (Figure \ref{fig:25-animal-pics}, third) and four-headed ram were found in Meroe West. The former is the only one found in the whole corpus, while the four-headed ram is the first in a non-elite context. An amulet featuring two joined frogs appears to be unfinished, as the two bases were not separated and each frog included its own suspension hole (Figure \ref{fig:25-animal-pics}, first). This is similar to the joined vulture amulet made of gold sheet from a Middle Napatan royal burial, which featured two suspension holes. An ibis amulet from Meroe West is very similar in design to the only other one in the corpus, from a contemporary royal grave---both feature a standing ibis holding a feather with its beak (Figure \ref{fig:25-animal-pics}, second).

Apart from those previously mentioned, other popular animal motifs included frogs, hawks, sows, apes, lizard, lions and cats, with more than 5 examples each. Other motifs were rarer, with less than 5 instances each, most with just one single example, such as falcons, geese, turtles, ibises, hippopotamus, jackals, monkeys and birds.


\begin{figure}[H]
	\centering
	\includegraphics[width=1\textwidth, height=10cm, keepaspectratio]{25-rest-motifs}
	\caption{Distribution of amulet nature, human, object and miscellaneous motifs in 25th Dynasty non-elite tombs.}
	\label{fig:25-rest-motifs}
\end{figure}


Amulets of nature, objects, humans and miscellaneous were included in the same graph as they were the least numerous (Figure \ref{fig:25-rest-motifs}). Once again, Meroe West featured the most distinct motifs (10), followed by Mirgissa (4), while Meroe South featured only three.


Amulets depicting humans and body parts were the most popular, in particular those of boys. They were made in mainly two designs and both were found in the same tombs. The first design features a rectangular-shaped head (Figure \ref{fig:25-boys-rectangular}), and the second displays a side lock of hair (Figure \ref{fig:25-boys-sidelock}). Both appear to be in the same half-seated position to Horus child. These amulets may indeed be representing this god and the different design is a Nubian adaptation without the thumb in mouth gesture. This hypothesis also applies to the twin boy amulets, the second most common in these categories. One design is very similar to the boy amulets, also featuring the side lock but not the hand in mouth motion, while the other design resembles "stick figures" and had no further details, but are depicted in the half-seated position (Figure \ref{fig:25-twin-boys}). This second design may also be a different representation of the same amulet, potentially of a double Horus child. Three tombs from Meroe West feature more than one pair of these amulets, while only one pair was found in a single Meroe South grave.


\begin{figure}[H]
	\centering
	\begin{minipage}{0.49\textwidth}
		\centering
		\includegraphics[width=\textwidth, height=4cm, keepaspectratio]{25-23-M-267}
	\end{minipage}
	\hspace{-20mm}%
	\begin{minipage}{0.49\textwidth}
		\includegraphics[width=\textwidth, height=4cm, keepaspectratio]{25-21.12103}
	\end{minipage}
	\caption{Boy amulets featuring a rectangular-shaped head (MFA 23-M-267, 21.12103).}
	\label{fig:25-boys-rectangular}
\end{figure}


\begin{figure}[H]
	\centering
	\begin{minipage}{0.45\textwidth}
		\centering
		\begin{minipage}{0.48\textwidth}
			\includegraphics[width=\textwidth, height=5cm, keepaspectratio]{25-47.3914}
		\end{minipage}%
		\begin{minipage}{0.48\textwidth}
			\includegraphics[width=\textwidth, height=5cm, keepaspectratio]{25-23-M-59}
		\end{minipage}
		\captionof{figure}{Amulet depicting boys with side lock (MFA 47.3914, 23-M-59).}
		\label{fig:25-boys-sidelock}
	\end{minipage}%
	\hspace{10pt}
	\begin{minipage}{0.45\textwidth}
		\centering
		\begin{minipage}{0.48\textwidth}
			\includegraphics[width=\textwidth, height=5cm, keepaspectratio]{25-23.368}
		\end{minipage}%
		\begin{minipage}{0.48\textwidth}
			\includegraphics[width=\textwidth, height=5cm, keepaspectratio, clip, trim={0cm 3cm 0cm 3cm}]{25-23.372}
		\end{minipage}
		\captionof{figure}{Amulets of twin boys with different designs (MFA 23.368, 23.372).}
		\label{fig:25-twin-boys}
	\end{minipage}
\end{figure}


Other human figures amulets include three depicting a naked female, two from Meroe West and one from Meroe South. The five amulets of a male with a double crown could be of an anthropomorphic deity or king, but the description did not specify any other features. The two motifs of human figures were unfortunately too obscure to be identified. One of them was of a plaque with an incised figure faded during baking (MFA 23-3-267). Even though, this amulet appears to have been damaged during the manufacture process it was still included in burial W 531 from Meroe West. Interestingly, this grave included only two amulets, the plaque with the human figure and a naked female. Eleven amulets depicted parts of the human body. Seven were of an incised face with a flower pattern, all from tomb W 643; and three depicted an eye and one a hand from Mirgissa. 

Object amulets include one offering table and one libation bucket from Meroe West and one sistrum from Mirgissa. The offering table amulet could be an affordable replacement for the real object, as the grave it was found in (W 671) did not include one. Papyrus columns were the only nature-related motif. Two came from Meroe West and one from Mirgissa. Miscellaneous motifs account for eight amulets. These include nine which were too fragmented or damaged to be identified. The remaining were four pierced golden nuggets and one amazonite pebble.


\paragraph{Amulets with multiple motifs}


\begin{figure}[H]
	\centering
	\includegraphics[width=1\textwidth]{25-2motifs}
	\caption{Distribution of amulet combining two motifs in 25th Dynasty non-elite tombs.}
	\label{fig:25-2motifs}
\end{figure}


In total, twelve amulets from 25th Dynasty non-elite assemblages feature a combination of two motifs. Figure \ref{fig:25-2motifs} displays the relation between the motifs and the quantity. For instance, the three plaque amulets of an ape also feature \textit{uraei} (Figure \ref{fig:25-2motifs-plaque}). The amulets are very similar in design, with the ape between the \textit{uraei} and either a snake or lotus flowers on top. Therefore, it is interesting that they were found in the same tomb, W 493 from Meroe West.


\begin{figure}[H]
	\centering
	\begin{minipage}{0.30\textwidth}
		\includegraphics[width=\textwidth, height=4cm, keepaspectratio]{25-23.762}
	\end{minipage}
	\hspace{-20mm}%
	\begin{minipage}{0.30\textwidth}
		\includegraphics[width=\textwidth, height=4cm, keepaspectratio]{25-23.763}
	\end{minipage}
	\hspace{-15mm}%
	\begin{minipage}{0.30\textwidth}
		\includegraphics[width=\textwidth, height=4cm, keepaspectratio]{25-2177}
	\end{minipage}
	\caption{Similar plaque amulets depicting an ape between two \textit{uraei} (MFA 23.762, 23.763, SNM 2177).}
	\label{fig:25-2motifs-plaque}
\end{figure}


\begin{wrapfigure}{l}{0.4\textwidth}
	\centering
	\includegraphics[width=1\linewidth, height=5cm, keepaspectratio, clip, trim={2cm 2.5cm 2cm 3cm}]{25-23.857}
	\caption{Plaque amulet of Bes and a winged goddess with the double crown (MFA 23.857).}
	\label{fig:25-2motifs-bes-isis}
\end{wrapfigure}


An openwork amulet features the god Bes holding two snakes with a winged goddess wearing the double crown on his side (Figure \ref{fig:25-2motifs-bes-isis}). Other two-motif amulets include a plaque with a giraffe and \textit{uraei}, and an amulet of a sun disc on top of a sun boat. Similarly to royal assemblages, the most common motif to include another was the Pataikos, with six amulets found in different graves. Among these, two feature Pataikos with the god Bes: one in which the gods are back-to-back and another with a bigger Pataikos holding a smaller Bes (Figure \ref{fig:25-2motifs-bes-pataikos}). The Pataikos is also featured in two other amulets with a hawk-headed deity on the reverse (Figure \ref{fig:25-2motifs-pataikos-hawk-head}). In fact, the Pataikos combined with the latter motif was feature in a double-sided amulet from a pre-25th Dynasty royal burial. Also similar to some royal amulets, two Pataikos' amulets were combined with winged leonine goddesses on the reverse.


\begin{figure}[H]
	\centering
	\begin{minipage}{0.30\textwidth}
		\includegraphics[width=\textwidth, height=5cm, keepaspectratio]{25-2181}
	\end{minipage}
	\hspace{-22mm}%
	\begin{minipage}{0.30\textwidth}
		\includegraphics[width=\textwidth, height=5cm, keepaspectratio]{25-2181-2}
	\end{minipage}
	\hspace{-18mm}%
	\begin{minipage}{0.30\textwidth}
		\includegraphics[width=\textwidth, height=5cm, keepaspectratio]{25-23.802}
	\end{minipage}
	\caption{Amulet of a back-to-back Bes and Pataikos (SNM 2181), and Pataikos holding a smaller Bes (MFA 23.802).}
	\label{fig:25-2motifs-bes-pataikos}
\end{figure}


\begin{figure}[H]
	\centering
	\begin{minipage}{0.20\textwidth}
		\centering
		\includegraphics[width=\textwidth, height=5cm, keepaspectratio]{25-23.771}
	\end{minipage}
	\hspace{-10mm}%
	\begin{minipage}{0.20\textwidth}
		\centering
		\includegraphics[width=\textwidth, height=5cm, keepaspectratio]{25-23.771-2}
	\end{minipage}
	\hspace{-8mm}%
	\begin{minipage}{0.20\textwidth}
		\centering
		\includegraphics[width=\textwidth, height=5cm, keepaspectratio]{25-23.828}
	\end{minipage}
	\hspace{-8mm}%
	\begin{minipage}{0.20\textwidth}
		\centering
		\includegraphics[width=\textwidth, height=5cm, keepaspectratio]{25-23.828-2}
	\end{minipage}
	\caption{Double-sided amulets of Pataikos and a hawk-headed deity (MFA 23.771, 23.828).}
	\label{fig:25-2motifs-pataikos-hawk-head}
\end{figure}



{\small
	\begin{xltabular}{\textwidth}{X X X X X}
		\caption{Distribution of amulet combining three motifs in 25th Dynasty non-elite tombs.}
		\label{tab:25-3motifs} \\
		\toprule
		\textbf{Site} & \textbf{Motif 1} & \textbf{Motif 2} & \textbf{Motif 3} & \textbf{Total} \\
		\midrule
		\endfirsthead
		\multicolumn{5}{c}{\textit{Continued from previous page}} \\
		\toprule
		\textbf{Site} & \textbf{Motif 1} & \textbf{Motif 2} & \textbf{Motif 3} & \textbf{Total} \\
		\midrule
		\endhead
		\bottomrule
		\endlastfoot
		Meroe West & four apes & hawk & ram & 1 \\
		\cmidrule(r){1-5}
		Meroe West & sun boat & winged \textit{uraeus} & scarab & 1 \\
		\cmidrule(r){1-5}
		Meroe West & hawk-headed deity & crocodile-headed deity & monkeys & 1 \\
		\cmidrule(r){1-5}
		Meroe South & Hathor & nb sign & winged \textit{uraeus} & 1 \\
		\cmidrule(r){1-5}
		Meroe South & four apes & hawk & ram & 1 \\
		\cmidrule(r){1-5}
		Mirgissa & cartouche & feathers & sun disc & 1 \\
	\end{xltabular}
}


A group of six amulets combined even more motifs (Table \ref{tab:25-3motifs}). They were distributed across five tombs, with W 609 containing two examples. With one exception, each unique design was a singular find. The exception is a pair of amulets depicting four apes supporting a ram, surrounded by hawks with outspread wings (Figure \ref{fig:25-3motifs-pics}, second). They were found in two separate burials from different sites (W 609 and S 127). The report even notes their similarity, raising the possibility that they were produced from the same mold \parencites[p.~429]{dunhamWestSouthCemeteries1963}.


Other amulets from this group include an openwork amulet showing deities---one with a hawk's head and another with a crocodile's head---each seated upon two monkeys holding a lotus flower together (Figure \ref{fig:25-3motifs-pics}, first). Another featured a scarab flanked by two winged \textit{uraei}, all on top of a sun boat (Figure \ref{fig:25-3motifs-pics}, third). These amulets were all from Meroe West, while Meroe South yielded a plaque depicting a Hathor head, also flanked by winged \textit{uraei}, surmounting the neb (“gold”) sign (Figure \ref{fig:25-3motifs-pics}, last). The presence of these amulets was not confined only to the Meroe cemeteries. A further example, a plaque showcasing a cartouche surrounded by feathers and a sun disc, was discovered in Tomb 22 at Mirgissa.


\begin{figure}[H]
	\centering
	\begin{minipage}{0.24\textwidth}
		\includegraphics[width=\textwidth, height=4cm, keepaspectratio]{25-23.801}
	\end{minipage}
	\hspace{-13mm}%
	\begin{minipage}{0.24\textwidth}
		\includegraphics[width=\textwidth, height=4cm, keepaspectratio]{25-23.819}
	\end{minipage}
	\hspace{-12mm}%
	\begin{minipage}{0.24\textwidth}
		\includegraphics[width=\textwidth, height=5cm, keepaspectratio]{25-23.820}
	\end{minipage}
	\hspace{-5mm}%
	\begin{minipage}{0.24\textwidth}
		\includegraphics[width=\textwidth, height=5cm, keepaspectratio, clip, trim={0cm 1.1cm 0cm 1.1cm}]{25-24.995}
	\end{minipage}
	\caption{Multi-motif amulets from 25th Dynasty non-elite tombs (MFA 23.801, 23.819, 23.820, 24.995).}
	\label{fig:25-3motifs-pics}
\end{figure}


\paragraph{Amulet materials}


Faience was, by far, the most common material, with 90\% of the total (Figure \ref{fig:25-faience}). As Meroe West contained the majority of amulets from this phase, naturally, it also boasted the most faience. Additionally, Meroe West also had the most diverse assemblage of amulet materials. In contrast, apart from faience, Mirgissa only had one other material (silver) and Meroe South, five (gold, a unidentified stone, silver, steatite and bone) (Figure \ref{fig:25-materials}). The only other material featured in all sites was silver; although, its presence was minimal at Meroe South and Mirgissa, with only one amulet each.

\begin{wrapfigure}{l}{0.5\textwidth}
	\centering
	\includegraphics[width=1\linewidth, height=10cm, keepaspectratio]{25-faience}
	\caption{Distribution of faience amulets in 25th Dynasty non-elite tombs.}
	\label{fig:25-faience}
\end{wrapfigure}


Gold, although the second most common material and sourced in Nubia, accounts for less than 2.6\% of the total amulets. They were found in the Meroe region, 36 in the West cemetery and three in the South. Ivory, also sourced from the Nubian territory, comprises only one amulet from Meroe West. Combined with the single ivory amulet form royal contexts, it seems this material was not preferred for the manufacture of amulets. Lapis lazuli, sourced from Afghanistan, was the material with the furthest source. The 36 lapis amulets were only found in Meroe West, but in different tombs, indicating a wider access to this material for this particular population.


Materials sourced from the Egyptian Eastern Desert---namely carnelian, jasper, serpentinite, steatite, porphyry, amethyst and, probably, silver---accounted for less then 3\% of the total amulets. Conversely, some stones could be obtained in quarries at the Lower Nubia region, such as rock crystal, amethyst and amazonite. The latter, as we have seen, could also be acquired from Ethiopia. However, there were only eight amulets made from these materials, less than 1\%. Bronze, although found in quarries in Northern Sudan as well as the Eastern Desert and Sinai \parencites[p.~935-938]{harrellArchaeologyGeologyAncient2024}, accounted for only three amulets---an ibis, a hawk and a boy. The ten amulets made of sandstone were all \textit{udjat} and came from the same tomb. In the report, the material is referred as quartzite, which is better applied to other types of sandstone. It is likely that this stone is the same used for two elite amulets from the Early-Middle Napatan phase.


New materials, previously not found in royal and elite tombs, include bone, amethyst, limestone, porphyry, glass and mother-of-pearl. Limestone likely came from the Nile Valley near the Egyptian city of Akhmin \parencites[p.~506]{harrellArchaeologyGeologyAncient2024} and it was the composition of only one amulet, the double Pataikos and hawk-headed deity from W 609. Mother-of-pearl, found in the Red Sea, and porphyry were used to produce two \textit{udjat} amulets. In fact, these three amulets are the only ones in the whole corpus to be made of these materials.


Conversely, the three amethyst amulets were found in different tombs and represented different motifs: an \textit{udjat}, a frog and a libation bucket. Amulets made of bone were present only at tomb S 177 and were all \textit{udjat}. A sole glass amulet of a lion was found in Meroe West. While glass production was technically related to the glazing of faience and steatite, it was a more complex and costly process. Thus, glass was mostly inaccessible to non-elite groups \parencites[p.~385]{harrellArchaeologyGeologyAncient2024}. Two hypothesis can be put forth regarding this amulet: the first is that it represents an exception for a non-elite context and the other is that the material was mislabeled. Nevertheless, the range of materials demonstrate the heterogeneity of access to objects and materials by non-elite groups, with some coming from further away, while others were readily available, like bone.


\begin{figure}[H]
	\centering
	\includegraphics[width=1\textwidth]{25-materials}
	\caption{Distribution of amulet materials (excluding faience) in 25th Dynasty non-elite tombs.}
	\label{fig:25-materials}
\end{figure}


\subsection{25th Dynasty-Early Napatan}


\subsubsection{Tomb structure}


\begin{figure}[H]
	\centering
	\includegraphics[width=1\textwidth]{25-EN-tombs}
	\caption{Distribution of super- and substructures in 25th Dynasty-Early Napatan non-elite tombs.}
	\label{fig:25-EN-tombs}
\end{figure}


\begin{figure}[H]
	\centering
	\includegraphics[width=1\textwidth]{25-EN-tomb-elements}
	\caption{Distribution of tomb features in 25th Dynasty-Early Napatan non-elite burials.}
	\label{fig:25-EN-tomb-elements}
\end{figure}


Non-elite tombs from the 25th Dynasty-Early Napatan phase were distributed across Meroe South (19), Meroe West (56), Missiminia (135), Qustul (32) and Sanam (1421) (Figure \ref{fig:25-EN-tombs}). All graves from the latter two sites were dated to this period. In total, there are 1674 burials. The structures were not much different from previous non-elite tombs, with the majority of substructures as pits. Although, two burials from Sanam were found in staircases of previous tombs and, therefore, were considered a different category. These burials could either be intrusive or belong to members of the same family who were interred afterwards. According to the reports, 72 graves from both cemeteries at Meroe lacked any superstructure. The remaining graves appear to have been covered by one, although its original form could not be determined.


These tombs exhibited a wider variety of architectural features than previous non-elite tombs. The distribution of these features, however, was highly uneven, with most occurring infrequently (Figure \ref{fig:25-EN-tomb-elements}). Trenches were the most widespread feature, present at most sites, with the exception of Qustul. These trenches may have functioned to accommodate burial bed legs or for the storage of grave goods. In contrast, bed holes were far rarer, found in just over 2\% of tombs and absent at Missiminia and Sanam.


The most common feature was the brick-lined pit, present in over 46\% of burials across Missiminia, Qustul, and Sanam. This was followed by brick roofs (13\%) and niches (8.1\%). Access to burial pits varied: at Missiminia, pits could be entered via stairs or slopes (with stairs being more popular), while at Sanam, if there was an entrance, it was always stairs.


The combination of features also differed by cemetery. The tombs at Meroe South and West had less elements, with 35 tombs containing only bed holes and trenches. Qustul was a little more diverse, with 19 graves featuring a combination of niches, bed holes and brick-lined pits. Sanam incorporated five different features, while Missiminia was the most heterogeneous, containing six different structural elements.


\subsubsection{Object types}


\begin{figure}
	\centering
	\includegraphics[width=1\textwidth]{25-EN-objs}
	\caption{Distribution of object types in 25th Dynasty-Early Napatan non-elite tombs.}
	\label{fig:25-EN-objs}
\end{figure}


25th Dynasty-Early Napatan assemblages feature 29 distinct object types, with Sanam being the most diverse (26 different objects) (Figure \ref{fig:25-EN-objs}). Meroe West, South and Qustul feature the least varied assemblages, with thirteen, ten and eight types, respectively. However, these object types were not distributed across many tombs as over half of the total did not contain any grave goods (54\%). Beads, adornments, vessels and scarabs are among the most widespread objects, with the remaining objects being very scarce in all sites. Only loose beads, adornments, vessels and tools are present in all four sites.


The variety and types of grave goods is comparable with 25th Dynasty royal tombs. This context, however, includes exclusive items, absent from previous phases and social contexts. They include incense burners, ostraca, palettes and seals, all from either Sanam and Missiminia, and were extremely rare. Additionally, most of the objects could have been used in daily life and, afterwards, buried with the deceased. Burial-exclusive objects, such as bead nets, canopics, offering tables and mummy trappings, were mainly found in Sanam and Missiminia. Burial containers, either preserved or traces of them were found in all sites, except Meroe South. Four tombs from Meroe West included a coffin, while the remaining three had burial beds. Traces of coffins were present in 30 tombs at Missiminia, while four tombs from Qustul suggested the use of burial beds. At Sanam, less than 3\% of burials contained remains of cartonnage which could indicate the use of a coffin, but only five of those tombs contained definite indications of a coffin.


\subsubsection{Amulets}


\paragraph{Amulet distribution}


This phase boasted the highest number of amulets, a total of 2682 distributed across 379 tombs (Table \ref{tab:25-EN-distro}). The Meroe cemeteries had the highest prevalence, with amulets found in over 30\% of tombs in each site. This is contrasted by Sanam, where, despite the large number of graves, only 23\% contained at least one amulet. The distribution was even lower at Qustul (17\%) and Missiminia (13\%). Despite the low prevalence of amulets in Sanam, the site contained the largest individual assemblages. For instance, tombs 1516 and 927 held over 200 amulets each.


{\small
	\begin{xltabular}{\textwidth}{X X X X X}
		\caption{Distribution of amulets in 25th Dynasty-Early Napatan non-elite tombs by site, including total counts of amulets and tombs, tombs containing amulets and amulet frequency (\%).}
		\label{tab:25-EN-distro} \\
		\toprule
		\textbf{Site} & \textbf{Total number of tombs} & \textbf{Total number of amulets} & \textbf{N. of tombs with amulets} & \textbf{\% of tombs with amulets} \\
		\midrule
		\endfirsthead
		\multicolumn{5}{c}{\textit{Continued from previous page}} \\
		\toprule
		\textbf{Site} & \textbf{Total number of tombs} & \textbf{Total number of amulets} & \textbf{N. of tombs with amulets} & \textbf{\% of tombs with amulets} \\
		\midrule
		\endhead
		\bottomrule
		\endlastfoot
		Meroe West & 57 & 147 & 21 & 37\% \\
		\cmidrule(r){1-5}
		Meroe South & 19 & 9 & 6 & 32\% \\
		\cmidrule(r){1-5}
		Missiminia & 135 & 123 & 17 & 13\% \\
		\cmidrule(r){1-5}
		Sanam & 1422 & 2390 & 330 & 23\% \\
		\cmidrule(r){1-5}
		Qustul & 32 & 13 & 5 & 17\% \\
		\midrule
		\textbf{Total} & \textbf{1665} & \textbf{2682} & \textbf{379} & -- \\
	\end{xltabular}
}


\paragraph{Amulet types}


\begin{figure}[H]
	\centering
	\includegraphics[width=1\textwidth]{25-EN-types}
	\caption{Distribution of amulet types in 25th Dynasty-Early Napatan non-elite tombs.}
	\label{fig:25-EN-types}
\end{figure}


Sanam and Meroe West were the only cemeteries where all amulet types were found, including the new category of human motifs (Figure \ref{fig:25-EN-types}). However, human amulets were not found in all sites, being absent from Meroe South and Missiminia. Four amulet types were present at Qustul, while only three were present at both Meroe South and Missiminia (animals, deities and symbols). At all sites, symbols were the most popular type of amulet. At Sanam, Qustul, and Missiminia, they are followed by deity amulets, but at Meroe West, the second most common are animals. In contrast, Meroe South, animal and deity motifs had only one instance each.


\paragraph{Amulet motifs}


\begin{figure}[H]
	\centering
	\includegraphics[width=1\textwidth]{25-EN-deity-motifs}
	\caption{Distribution of amulet deity motifs in 25th Dynasty-Early Napatan non-elite tombs.}
	\label{fig:25-EN-deity-motifs}
\end{figure}


A total of 24 distinct deities are depicted in 25th Dynasty-Early Napatan amulets (Figure \ref{fig:25-EN-deity-motifs}). Apart from those, there are 29 amulets of unidentified deities, including one with a group of deities. Geographically, they were found in all sites, with Sanam accounting for 23 motifs, of which thirteen were only found there. Meroe West comprised ten different motifs, without any exclusive to this site, while Meroe South only included a fragment of a lion-headed deity. Qustul was home to two deity amulets, one of Bes and another of Hathor. At Missiminia, there were amulets of six deities: Pataikos, Bes, Isis, Hathor, Isis and Horus, and Taweret. In contrast to previous phases, amulets of Sekhmet were the most popular, followed by Pataikos, Bes, Isis and Hathor. This is also the only phase to include amulets of the deities Khonsu, Neith, and Min, all from Sanam.


\begin{figure}[H]
	\centering
	\includegraphics[width=1\textwidth, height=4cm, keepaspectratio]{25-EN-XXVI.27-31-1516}
	\caption{Plaque amulets of Hapi, the Nile god, from Sanam \parencites[pl.~XXVI.27-31]{griffithOxfordExcavationsNubia1923}}
	\label{fig:25-EN-hapi-plaques}
\end{figure}


Regarding specific motifs, there are eight plaque amulets from Sanam of Hapi, the Nile god, which are the first ones found in any burials since one instance from a 25th Dynasty royal grave (Figure \ref{fig:25-EN-hapi-plaques}). Similarly, two winged lion-headed goddess, also from Sanam, are the only ones from a non-royal context (Figure \ref{fig:25-EN-winged-lioness}). They are depicted in different positions, similar to the royal ones, one sideways and another facing frontward, with a prominent belly and tights. Both also have the sun disc and \textit{uraeus} on top of their heads.

\begin{wrapfigure}{l}{0.5\textwidth}
	\centering
	\includegraphics[width=0.49\textwidth, height=5cm, keepaspectratio]{25-EN-XXVI32}
	\includegraphics[width=0.49\textwidth, height=5cm, keepaspectratio]{25-EN-LV13}
	\caption{Winged lion-headed goddess amulets, from Sanam \parencites[pl.~XXVI.32, LV.13]{griffithOxfordExcavationsNubia1923}}
	\label{fig:25-EN-winged-lioness}
\end{wrapfigure}

One particular motif is puzzling as it depicts a winged goddess with a child, which could potentially represent the standard Isis and Horus motif. The first picture in Figure \ref{fig:25-EN-isis-horus} show the classic design of a seated Isis, crowned with a sun disc and \textit{uraeus}, holding Horus on her lap. The second image, from Missiminia, displays the intriguing amulet of a winged goddess wearing the Hathor crown with a child either seated by her feet or on her lap. Vila raises the possibility that it represents Hathor presenting a child Khonsu \parencites[p. 53]{vilaNécropoleMissiminiaSépultures1980}. The last amulet from \ref{fig:25-EN-isis-horus} is also intriguing. It depicts two figures side-by-side, the one on the left also sporting the Hathor/Isis crown. They appear to be staring and embracing each other. This motif is particularly intriguing as nothing like it was found. Additionally, Griffith describes it as “Isis nursing Horus, barbaric” \parencites[p.~135]{griffithOxfordExcavationsNubia1923}. Unfortunately, I was not able to find the current location of this amulet in order to analyze a better picture. This amulet is also similar to some erotic amulets, where two human figures are next to each other and one seems to be kissing the other. It is possible that it may represent an intimate scene between Isis and Osiris. Additionally, the incisions on their bodies could be an indication of mummification, reinforcing the Osiris interpretation. In any case, both amulets are very unique and their design look very Nubian in style.


\begin{figure}[H]
	\centering
	\begin{minipage}{0.30\textwidth}
		\includegraphics[width=\textwidth, height=4cm, keepaspectratio]{25-EN-23.337}
	\end{minipage}
	\hspace{-25mm}%
	\begin{minipage}{0.30\textwidth}
		\includegraphics[width=\textwidth, height=4cm, keepaspectratio]{25-EN-37-5}
	\end{minipage}
	\hspace{-20mm}%
	\begin{minipage}{0.30\textwidth}
		\includegraphics[width=\textwidth, height=4cm, keepaspectratio]{25-EN-LVI.15}
	\end{minipage}
	\caption{Amulets of Isis and Horus in different designs, from Meroe West (MFA 23.337), Missiminia \parencites[p.~52, fig. 37, 5]{vilaNécropoleMissiminiaSépultures1980}, and Sanam \parencites[pl.~LVI.15]{griffithOxfordExcavationsNubia1923}}
	\label{fig:25-EN-isis-horus}
\end{figure}


Other individual amulets of deities include one of a hawk-headed dwarf wearing the \textit{hmhm} crown, which could be a representation of the god Bes. This motif is similar to the ram-headed dwarf amulet form a 25th Dynasty royal burial, although both sport different crowns and animal heads. A amulet of Bes from Missiminia is represented in a novel Nubian style found especially in post-New Kingdom non-elite burials from Buhen, Amara West, and the Faras and Kawa/Napata regions \parencites{petacchiLocalAspectsCult2014}. Similarly, an amulet of a Hathor head from Qustul also featured a very unique representation (Figure \ref{fig:25-EN-deity-pics})


\begin{figure}[H]
	\centering
	\begin{minipage}{0.3\textwidth}
		\includegraphics[width=\textwidth, height=4cm, keepaspectratio]{25-EN-LV9}
	\end{minipage}
	\hspace{-30mm}%
	\begin{minipage}{0.3\textwidth}
		\includegraphics[width=\textwidth, height=4cm, keepaspectratio]{25-EN-61-4}
	\end{minipage}
	\hspace{-30mm}%
	\begin{minipage}{0.3\textwidth}
		\includegraphics[width=\textwidth, height=4cm, keepaspectratio]{25-EN-VB25-2}
	\end{minipage}
	\caption{Amulets depicting a hawk-headed dwarf, Bes and Hathor \parencites[pl.~LV.9]{griffithOxfordExcavationsNubia1923}[p.~69, fig. 61, 4]{vilaNécropoleMissiminiaSépultures1980}[p.~74, fig. 25, c]{williamsExcavationsAbuSimbel1990}}
	\label{fig:25-EN-deity-pics}
\end{figure}


\begin{figure}[H]
	\centering
	\includegraphics[width=1\textwidth]{25-EN-udjat}
	\caption{Distribution of amulet \textit{udjat} and quadruple \textit{udjat} motifs in 25th Dynasty-Early Napatan non-elite tombs.}
	\label{fig:25-EN-udjat}
\end{figure}


\begin{wrapfigure}{l}{0.3\textwidth}
	\centering
	\includegraphics[width=1\linewidth, height=5cm, keepaspectratio]{25-EN-F1940.12.65}
	\caption{Examples of a \textit{udjat} and quadruple \textit{udjat} amulets from Sanam (Leiden F 1940/12.65).}
	\label{fig:25-EN-udjat-pic}
\end{wrapfigure}


Amulets of the \textit{udjat} and quadruple \textit{udjat} were included in a separate graph, as they were among the most numerous (Figure \ref{fig:25-EN-udjat}). For the latter motif, four \textit{udjat} were combined in order to form a single amulet (Figure \ref{fig:25-EN-udjat-pic}). \textit{Udjat} amulets were found in all five sites, although in different quantities. Percentage wise, the udjat comprises 35\% of all amulets from this phase. Qustul and Meroe South contained the least amount of this motif (four and seven, respectively), while Missiminia accounted for 28, Meroe West for 46 and Sanam with the highest amount, 845.


Meroe South was the only site without any amulets of the quadruple \textit{udjat}---although Meroe West contained only one and Qustul, two. This absence may indicate regional variations in amulet motifs, which will be explored in a later section. There is a total of 122 examples of the quadruple \textit{udjat}, which were rarer in burials dated solely to the 25th Dynasty, which may indicate that the motif was widespread in both phases, 25th Dynasty and Early Napatan.


\begin{figure}[H]
	\centering
	\includegraphics[width=1\textwidth]{25-EN-symbol-motifs}
	\caption{Distribution of amulet symbol motifs (excluding \textit{udjat} and quadruple \textit{udjat}) in 25th Dynasty-Early Napatan non-elite tombs.}
	\label{fig:25-EN-symbol-motifs}
\end{figure}


The remaining symbol amulets included several motifs previously restricted to royal or elite burials (Figure \ref{fig:25-EN-symbol-motifs}). These include the heart, \textit{djed}, \textit{menat} and \textit{nefer} amulets. Furthermore, winged scarab and the ram-headed, and human-headed sphinxes motifs were previously only known from 25th Dynasty royal graves. Amulets of the \textit{ankh} and \textit{sa} signs, which had been confined to Early Napatan elite tombs, are also present. Additionally, this phase introduced unique motifs such as the horned lunar disc, lunar crescent, and sun disc. Among these, the most common motif was the heart, comprising of only eight examples though.


The distribution of these amulets across sites varied. Sanam was the most diverse, with eleven motifs, seven of which were exclusive to the site. Missiminia included five (three exclusive), while Meroe West accounted for three motifs, none of which were unique to the site. No other symbol amulets---apart from the \textit{udjat} and quadruple \textit{udjat}---were found in Qustul or Meroe South.


Amulets depicting animals are much more diverse than previous non-elite assemblages as well as those of royalty and elite (Figure \ref{fig:25-EN-animal-motifs}). All sites from this phase included at least one animal amulet, as is the case with Qustul, where only one crocodile amulet was found, and Meroe South, with one ram amulet. Sanam contains not only the most diverse assemblage of animal amulets, 26 out of the 29 total motifs, but also the biggest amounts. Furthermore, Sanam was home to 15 exclusive motifs. Meroe West comprised 11 motifs, with a single exclusive one, a crocodile with a hawk's head. Missiminia included six motifs, with birds being exclusive to the site.


Frogs are the most widespread animal motif, present in Sanam, Meroe West and Missiminia. Four amulets depicting two frogs, joined either side-by-side or behind, were found at Sanam (Figure \ref{fig:25-EN-frogs}). They were followed by cats, flies, sows, rams, scorpions and monkeys, all with more than 20 instances each. It is notable that these motifs mostly depict regional animals and not necessarily associated with an Egyptian deity, apart from cats and rams. Amulets of sows were oftentimes depicted with piglets. Furthermore, two hippopotamus amulets were found at Meroe West and Sanam and their design is similar to the sow ones, featuring both quadrupeds in a walking pose, with their nose on the ground, as if searching for food (Figure \ref{fig:25-EN-sows}). In fact, the only other hippo amulet in the corpus was found in a non-elite 25th Dynasty burial. Other exclusive non-elite amulets to this phase also depict regional animals, such as bulls, crocodiles, the double frogs, ducks, fishes, hares, hedgehogs, pigs, scorpions and snakes, comprising 16\% of all animal amulets from this particular context.


\begin{figure}[H]
	\centering
	\includegraphics[width=1\textwidth]{25-EN-animal-motifs}
	\caption{Distribution of amulet animal motifs in 25th Dynasty-Early Napatan non-elite tombs.}
	\label{fig:25-EN-animal-motifs}
\end{figure}

\begin{figure}[H]
	\centering
	\begin{minipage}{0.30\textwidth}
		\includegraphics[width=1\textwidth, height=3cm, keepaspectratio]{25-EN-frogs}
		\caption{Double frog amulets from Sanam \parencites[pl.~LVII.12, 15]{griffithOxfordExcavationsNubia1923}}
		\label{fig:25-EN-frogs}
	\end{minipage}
	\hspace{15pt}
	\begin{minipage}{0.45\textwidth}
		\includegraphics[width=0.49\textwidth, height=4cm, keepaspectratio, clip, trim={10cm 10cm 10cm 10cm}]{25-EN-23.358}
		\includegraphics[width=0.49\textwidth, height=4cm, keepaspectratio]{25-EN-LVI.40}
		\captionof{figure}{Amulets of a hippopotamus from Meroe West (MFA 23.358) and a sow with piglets from Sanam \parencites[pl.~LVI.40]{griffithOxfordExcavationsNubia1923}}
		\label{fig:25-EN-sows}
	\end{minipage}
\end{figure}


Scarabs as amulets make their first appearance in this assemblage. These motifs could be depicted in different forms. For instance, Sanam included scarabs with spread legs and one with the legs tucket in and a ram's head (Figure \ref{fig:25-EN-scarabs}). One of the ram amulets found at Missiminia has a unique design, different from those at royal tombs, although it does include an \textit{uraeus} (Figure \ref{fig:25-EN-animal-pics}, second). Three four-headed rams were found at Sanam, which were previously found in royal tombs of the Middle Napatan phase and non-elite ones from the 25th Dynasty.

Some motifs are particular noteworthy as they combined different animal features. For instance, a crocodile with a hawk's head, lion-headed flies (Figure \ref{fig:25-EN-animal-pics}, first and third) and the scarab with a ram's head. Apart from the lion-headed flies, the remaining motifs have no parallels in the corpus.


\begin{figure}[H]
	\centering
	\begin{minipage}{0.25\textwidth}
		\centering
		\includegraphics[width=\linewidth, height=3cm, keepaspectratio]{25-EN-scarabs}
		\caption{Scarab amulets from Sanam, one with spread legs and another with a ram's head \parencite[pl.~LVII.30, 32]{griffithOxfordExcavationsNubia1923}}
		\label{fig:25-EN-scarabs}
	\end{minipage}%
	\hspace{20pt}
	\begin{minipage}{0.50\textwidth}
		\centering
		\begin{minipage}{0.40\textwidth}
			\includegraphics[width=\linewidth, height=4cm, keepaspectratio]{25-EN-LVII.21,22}
		\end{minipage}%
		\begin{minipage}{0.30\textwidth}
			\includegraphics[width=\linewidth, height=4cm, keepaspectratio]{25-EN-211-1}
		\end{minipage}%
		\begin{minipage}{0.55\textwidth}
			\includegraphics[width=\linewidth, height=4cm, keepaspectratio, clip, trim={0.5cm 0.5cm 0.5cm 0.5cm}]{25-EN-47.3023}
		\end{minipage}
		\caption{Amulets of lion-headed flies, a ram's head with \textit{uraeus}, and a hawk-headed crocodile (\cites[pl.~LVII.21, 22]{griffithOxfordExcavationsNubia1923}[p. 87, fig. 87, 2]{vilaNécropoleMissiminiaSépultures1980}, MFA 47.3023)}
		\label{fig:25-EN-animal-pics}
	\end{minipage}
\end{figure}


\begin{figure}[H]
	\centering
	\includegraphics[width=1\textwidth]{25-EN-rest-motifs}
	\caption{Distribution of amulet nature, object, human and other motifs in 25th Dynasty-Early Napatan non-elite tombs.}
	\label{fig:25-EN-rest-motifs}
\end{figure}


Among the remaining amulet categories, Sanam stands out, once again, with the greatest diversity, featuring 17 different ones---13 of which are exclusive to the site (Figure \ref{fig:25-EN-rest-motifs}). In contrast, Meroe West yielded eight different motifs (five exclusive) and Qustul, which included only one eye amulet. No amulets from these categories were found at Missiminia or Meroe South.


\begin{figure}[H]
	\centering
	\begin{minipage}{0.25\textwidth}
		\includegraphics[width=\textwidth, height=4cm, keepaspectratio]{25-EN-E.05716.4a}
	\end{minipage}
	\hspace{-18mm}%
	\begin{minipage}{0.25\textwidth}
		\includegraphics[width=\textwidth, height=4cm, keepaspectratio]{25-EN-E.131.1921}
	\end{minipage}
	\hspace{-18mm}%
	\begin{minipage}{0.25\textwidth}
		\includegraphics[width=\textwidth, height=4cm, keepaspectratio]{25-EN-LVIII.22}
	\end{minipage}
	\hspace{-28mm}%
	\begin{minipage}{0.25\textwidth}
		\includegraphics[width=\textwidth, height=4cm, keepaspectratio]{25-EN-24.831}
	\end{minipage}
	\caption{Amulets from Sanam: papyrus column, sistrum and harpoon (Brüssels E.05716.4a, Fitzwilliam E.131.1921, \parencites[pl.~LVIII.22]{griffithOxfordExcavationsNubia1923}, and MFA 24.831)}
	\label{fig:25-EN-rest-pics}
\end{figure}


Nature motifs feature papyrus columns, lotus flowers, a pomegranate and an unidentified flower. Papyrus columns and lotus flowers are among the most prevalent motifs, with 66 instances. Among these, Sanam included all examples of nature motifs, with Meroe West containing only one of the papyrus columns. The same pattern is observed with object amulets. Sistrum, offering table, harpoon, vase, libation bucket, and jar motifs were only featured in Sanam, while all dovetails were found at Meroe West (some examples in Figure \ref{fig:25-EN-rest-pics}).



\begin{figure}[H]
	\centering
		\includegraphics[width=\textwidth, height=4cm, keepaspectratio, clip, trim={1cm 3.3cm 1cm 1cm}]{25-EN-AN1921.75.21.35}
		\caption{Amulets of eyes and double eyes from Sanam (Ashmolean AN1921.75.21.35).}
		\label{fig:25-EN-eyes}
\end{figure}


\begin{figure}[H]
	\centering
		\includegraphics[width=\textwidth, height=4cm, keepaspectratio]{25-EN-humans}
		\caption{Amulets of a man on a horse and a nude female, from Sanam  \parencites[pl.~XXIX.3, LVI.14]{griffithOxfordExcavationsNubia1923}}
		\label{fig:25-EN-humans}
\end{figure}


Human figures and body parts were more distributed between the sites. Most eye and hand motifs, and all instances of double eyes, naked female, a human with a sun disc (possibly an unidentified deity) and man on a horse were found at Sanam (Figures \ref{fig:25-EN-humans}, \ref{fig:25-EN-eyes}). At Meroe West, there were two amulets of hands and all amulets of boys and twin boys. Only one eye amulet came from Qustul. The man on horse motif is the only one in the corpus and the naked female is very Nubian in design, similar to figurines, with incised decoration, and prominent belly and thighs. Among the miscellaneous motifs are an amulet case stuffed with organic material and a golden nugget, both from Meroe West.


\paragraph{Amulets with multiple motifs}


\begin{figure}[H]
	\centering
	\includegraphics[width=1\textwidth]{25-EN-2motifs}
	\caption{Distribution of amulet combining two motifs in 25th Dynasty-Early Napatan non-elite tombs.}
	\label{fig:25-EN-2motifs}
\end{figure}


A total of 55 amulets from Sanam, Missiminia, Qustul and Meroe West feature two motifs, fourteen of which are plaques. The most popular is the combination of \textit{udjat} with \textit{ankh}, which appears in 21 amulets (Figure \ref{fig:25-EN-2motifs}). The remaining combinations are far rarer, with the majority being unique.


\begin{figure}[H]
	\centering
	\begin{minipage}{0.30\textwidth}
		\includegraphics[width=\textwidth, height=4cm, keepaspectratio]{25-EN-24.781}
	\end{minipage}
	\hspace{-25mm}%
	\begin{minipage}{0.30\textwidth}
		\includegraphics[width=\textwidth, height=4cm, keepaspectratio]{25-EN-24.781-2}
	\end{minipage}
	\hspace{-20mm}%
	\begin{minipage}{0.30\textwidth}
		\includegraphics[width=\textwidth, height=4cm, keepaspectratio]{25-EN-320-13}
	\end{minipage}
	\caption{Amulet of a Pataikos crowned by a winged scarab, with an incised jackal on the reverse; and a Bes head on top of Amun, Mut and Khonsu (MFA 24.781, \cites[p.~128, fig. 139, 2]{vilaNécropoleMissiminiaSépultures1980})}
	\label{fig:25-EN-2motifs-pics}
\end{figure}


Similarly to some royal and non-elite 25th Dynasty amulets, three examples depict a winged lion-headed goddess on the back of a Pataikos. A particular amulet even features two Pataikos back-to-back. Animal motifs are frequent, appearing in combinations, such as a hippopotamus with a lion and a crocodile; an \textit{udjat} with a double ram and apes; and a Pataikos with a jackal (Figure \ref{fig:25-EN-2motifs-pics}, first). A ram-headed sphinx is also paired with an \textit{uraeus}.


Deities are well represented in these amulets, with a variety including Amun, Mut, Khonsu, Hathor, Horus, Onuris, Pataikos, Ra-Horakhty, Sekhmet, Thoth, Bes and Tefnut. One particularly unique amulet depicts the Theban triad (Amun, Mut and Khonsu) positioned beneath the head of Bes, or, perhaps, being embraced by the god (Figure \ref{fig:25-EN-2motifs-pics}, second). Finally, symbols of natural elements with symbolic meanings like lunar crescents, feathers, lotus and sun discs were also paired with deities.


Amulets featuring three motifs were exclusive to Sanam and Meroe West, totaling 27 instances (Table \ref{tab:25-EN-3motifs}). The most common combination was the sun boat with an \textit{uraeus} and \textit{hmhm} crown, of which the 15 total instances were found in a single grave from Sanam. The \textit{udjat} appeared in the most varied combinations, paired with elements such as a basket and papyrus; cat and vulture wings (Figure \ref{fig:25-3motifs-pics}, top left); arms and the \textit{ankh}; and with an \textit{uraeus} and ram. Another amulet has the sun disc between two \textit{uraeus} on top of lotus flowers (Figure \ref{fig:25-3motifs-pics}, top right).


{\small
	\begin{xltabular}{\textwidth}{X X X X X}
		\caption{Distribution of amulet combining three motifs in 25th Dynasty-Early Napatan non-elite tombs.}
		\label{tab:25-EN-3motifs} \\
		\toprule
		\textbf{Site} & \textbf{Motif 1} & \textbf{Motif 2} & \textbf{Motif 3} & \textbf{Total} \\
		\midrule
		\endfirsthead
		\multicolumn{5}{c}{\textit{Continued from previous page}} \\
		\toprule
		\textbf{Site} & \textbf{Motif 1} & \textbf{Motif 2} & \textbf{Motif 3} & \textbf{Total} \\
		\midrule
		\endhead
		\bottomrule
		\endlastfoot
		Sanam & ram-headed sphinx & sun boat & jackal & 1 \\
		\cmidrule(r){1-5}
		Sanam & horus & scorpion & crocodile & 1 \\
		\cmidrule(r){1-5}
		Sanam & sun boat & \textit{uraeus} & hmhm crown & 15 \\
		\cmidrule(r){1-5}
		Sanam & fish & \textit{uraeus} & crocodile & 1 \\
		\cmidrule(r){1-5}
		Sanam & monkey & \textit{uraeus} & crocodile & 1 \\
		\cmidrule(r){1-5}
		Sanam & pillar & sun disc & feathers & 1 \\
		\cmidrule(r){1-5}
		Sanam & hathor & \textit{uraeus} & winged scarab & 1 \\
		\cmidrule(r){1-5}
		Sanam & \textit{udjat} & cat & vulture wings & 1 \\
		\cmidrule(r){1-5}
		Sanam & \textit{udjat} & \textit{uraeus} & ram & 1 \\
		\cmidrule(r){1-5}
		Sanam & \textit{udjat} & basket & papyrus & 2 \\
		\cmidrule(r){1-5}
		Sanam & \textit{udjat} & arms & ankh & 1 \\
		\cmidrule(r){1-5}
		Sanam & sun disc & \textit{uraeus} & lotus & 1 \\
		\cmidrule(r){1-5}
		Qustul & hippo & jackal & crocodile & 1 \\
	\end{xltabular}
}


The \textit{uraeus} was the most frequently combined motif, often paired with animals and nature symbols such as monkeys, fishes and crocodiles. This thematic connection is also seen in a plaque featuring only animals: a hippopotamus, jackal and crocodile. A particular plaque shows a ram-headed sphinx on the front and a hieroglyphic inscription on the back (\ref{fig:25-3motifs-pics}, a, b). Deities were also represented, including a plaque of Horus standing on a crocodile with two scorpions on the reverse (\ref{fig:25-3motifs-pics}, c, d), and another depicting Hathor with an \textit{uraeus} and a winged scarab. Finally, one unique amulet was described as a pillar crowned with a disc and feathers.


\begin{figure}[H]
	\centering

	\begin{subfigure}{0.45\textwidth}
		\centering
		\includegraphics[width=\linewidth, height=3.5cm, keepaspectratio]{25-EN-E.05714.4}
	\end{subfigure}%
	\hspace{5pt}
	\begin{subfigure}{0.45\textwidth}
		\centering
		\includegraphics[width=\linewidth, height=3.5cm, keepaspectratio]{25-EN-E.05716.24}
	\end{subfigure}
	
	\vspace{8pt}

	\begin{subfigure}{0.45\textwidth}
		\centering
		\includegraphics[width=\linewidth, height=4cm, keepaspectratio]{25-EN-E.05701.7}
		\caption{Front}
	\end{subfigure}%
	\hspace{5pt}
	\begin{subfigure}{0.45\textwidth}
		\centering
		\includegraphics[width=\linewidth, height=4cm, keepaspectratio]{25-EN-E.05701.7-2}
		\caption{Back}
	\end{subfigure}
	
	\vspace{8pt}

	\begin{subfigure}{0.45\textwidth}
		\centering
		\includegraphics[width=\linewidth, height=4cm, keepaspectratio]{25-EN-F1940-11.10}
		\caption{Front}
	\end{subfigure}%
	\hspace{5pt}
	\begin{subfigure}{0.45\textwidth}
		\centering
		\includegraphics[width=\linewidth, height=3.5cm, keepaspectratio]{25-EN-F1940-11.10-2}
		\caption{Back}
	\end{subfigure}
	
	\caption{Amulets with at least three motifs (front and back views) (Brüssels E.05714.4, E.05716.24, E.05701.7, Leiden F 1940/11.10)}
	\label{fig:25-EN-3motifs-pics}
\end{figure}


\paragraph{Amulet materials}


\begin{figure}[H]
	\centering
	\includegraphics[width=1\textwidth]{25-EN-faience}
	\caption{Distribution of faience amulets in 25th Dynasty-Early Napatan non-elite tombs.}
	\label{fig:25-EN-faience}
\end{figure}


\begin{figure}[H]
	\centering
	\includegraphics[width=1\textwidth]{25-EN-materials}
	\caption{Distribution of amulet materials (excl. faience) in 25th Dynasty-Early Napatan non-elite tombs.}
	\label{fig:25-EN-materials}
\end{figure}


Faience was, once again, the most widespread material for amulets, comprising over 82\% of the total assemblage and present at all sites (Figure \ref{fig:25-EN-faience}). Sanam yielded the most amulets, 2007, which is proportionate to the number of amulets. It was followed by 99 instances from Meroe West, 97 from Missiminia, eight from Qustul and seven from Meroe South. The remaining amulets were made from 29 different materials, with Sanam exhibiting the greatest diversity by far (26) (Figure \ref{fig:25-EN-materials}). As with motifs, Sanam also had the most exclusive materials, with 13 not found in the other sites. Meroe West included 12 distinct materials, though none were unique to the site. Missiminia's amulets were composed of six materials, including malachite as an exclusive. Qustul featured steatite, lapis lazuli and jasper, while only steatite and ivory were present at Meroe South.


The variety of materials showcases extensive trade networks. Of the 29 materials, 13 were imported from beyond the Nubian territory. These could come from relatively nearby sources like the Egyptian Eastern Desert, Fayum or Sinai (e.g., amazonite, rock crystal, steatite, malachite, quartz, basalt, amethyst, chalcedony, haematite, jasper, travertine) \parencites[p.~505-507, 758-760]{harrellArchaeologyGeologyAncient2024}. However, they made up less than 7\% of the assemblage.


On the other hand, lapis lazuli and obsidian, coming from much further lands, accounted for 2.4\% of amulets \parencites[p.~759]{harrellArchaeologyGeologyAncient2024}. Obsidian was unique to this non-elite context and was found only at Sanam. Eight materials, such as silver, bronze, agate, carnelian, zeolite, granite and onyx, are more ambiguous,as they could be sourced from either Nubia or Egypt \parencites{andrewNoteZeoliteArtefacts1951}[p.~505, 759, 875-876, 929]{harrellArchaeologyGeologyAncient2024}. This phase is also the third to include ivory amulets, four in total, one from Meroe South and three from Meroe West. Gold amulets were not very popular, making up only 0.7\% of the total. Notably, glass---a more costly material than faience---was not only more common in this phase, with 32 instances versus one in the previous phase.


Despite the abundance of imported and local materials, Sanam also yielded five amulets made of readily available resources such as shell, pottery and bone. Additionally, Sanam's remarkable diversity of materials for amulets could be explained by the sizable assemblage, as well as its key position as a trading post, facilitating access to both local and long-distance exchange networks.


\subsection{25th Dynasty-Middle Napatan}


\subsubsection{Tomb structure}


The 159 non-elite tombs dated to the 25th Dynasty-Middle Napatan range were attested only at Meroe West. These graves have the broadest span of time, encompassing three phases. As most of these tombs were found completely plundered, it is possible that they were dated to this range by Reisner/Dunham based on their position to other tombs rather than through material culture. Unfortunately, dating these tombs relying solely on the report is not possible, and, therefore, we depend on the author's assessment. Furthermore, all graves were simple pits without superstructure, and the majority did not include any further features, with the exception of two tombs: W 499, which was lined with bricks, and W 794, which included roofing.


\subsubsection{Object types}


Although 25th Dynasty-Middle Napatan non-elite assemblages generally follow the Napatan pattern we have seen so far, they are smaller in scale than their elite and royal counterparts (Figure \ref{fig:25-MN-objs}). This reduction in the quantity of grave goods likely resulted from either looting or intentional deposition of fewer objects. Additionally, the scarcity of grave goods is underscored by the fact that 88 graves (over 55\% of the total) were completely empty.


Loose beads continue to be the most popular object type, with a total of 1932 instances. Ceramic and non-ceramic vessels, are less common than in other social groups. Other types of objects were extremely rare, with less than 10 instances each distributed in the 68 furnished graves.


Objects made for the burial, such as canopics and containers, are absent, with the exception of bead nets and a single offering table. Bead nets were found in 13\% of the burials and, since they were constructed from readily available beads, their production was likely unrestricted. Furthermore, many of the loose beads found in other graves may have been part of disintegrated nets. The unique offering table was decorated with cakes and Hes-vases, but with no trace of inscription. Like the unique offering table in 25th Dynasty burials, the objects were not common possessions of the non-elite and were most likely acquired by individuals who had social or trade connections.

\begin{wrapfigure}{l}{0.5\textwidth}
	\centering
	\includegraphics[width=1\linewidth, height=8cm, keepaspectratio]{25-MN-objs}
	\caption{Distribution of object types in 25th Dynasty-Middle Napatan non-elite tombs.}
	\label{fig:25-MN-objs}
\end{wrapfigure}

However, non-elite assemblages of this phase do include one object category which was completely absent in elite burials and previous non-elite ones, but present in royal assemblages from all phases: offering vessels. This may be due to two possibilities: the first being a difference in terminology, as Reisner/Dunham explicitly categorize certain objects as “offering vessels”, ascribing a function to them, whereas reports from other sites do not use this exact term. Another possibility is that this object type is, in fact, absent from previous non-elite assemblages, as the reports from at Meroe West and South, composed by Reisner/Dunham, did not include this type of object. If the absence in earlier tombs is genuine, it suggests that either non-elite access and use of this object expanded during this specific period, and that the tombs containing them should be dated later than the 25th Dynasty.


\subsubsection{Amulets}


\paragraph{Amulet distribution}


Amulets deposition decreased in this phase. The total 192 amulets were distributed in only ten of the total 159 tombs, only 6\%. More precisely, two tombs account for more than 89\% of the amulets, W 734 and W 654, which contained 141 and 32 amulets, respectively. The remaining six tombs with amulets had less than five instances each.


\paragraph{Amulet types}


The 25th Dynasty-Middle Napatan amulets depicted six different categories (Figure \ref{fig:25-MN-types}). Amulets representing human body parts were the most popular, with 43\% of the total. Deities were the second most common type, comprising of 39\% of all amulets, whereas symbols made up 22\%. The remaining categories were much rarer, with miscellaneous motifs making up 4\%, animals 3\% and a sole nature-related motif.

\begin{figure}[H]
	\centering
	\includegraphics[width=1\textwidth]{25-MN-types}
	\caption{Distribution of amulet types in 25th Dynasty-Middle Napatan non-elite tombs.}
	\label{fig:25-MN-types}
\end{figure}



\paragraph{Amulet motifs}


\begin{figure}[H]
	\centering
	\includegraphics[width=1\textwidth]{25-MN-motifs}
	\caption{Distribution of amulet motifs in 25th Dynasty-Middle Napatan non-elite tombs.}
	\label{fig:25-MN-motifs}
\end{figure}


Eye amulets were the only motif in the human category and, as we have seen, were the most popular (Figure \ref{fig:25-MN-motifs}). Additionally, they were all strung in the same necklace. They were followed by amulets of the gods Pataikos and Bes, encompassing 38\%. Similar to the eye amulets, all of the Bes amulets were strung together into a necklace, along with a single sow amulet (Figure \ref{fig:25-MN-pics}). The \textit{udjat}, although one of the most popular motif in earlier phases, made up only 11\% in this phase.


Among the animals are four amulets of sows, one of a monkey and another of a scarab. Additionally to this phase non-elite graves, sow and scarab amulets were found in other contexts, non-elite and elite assemblages from the 25th Dynasty to the Early Napatan phases. As they were not included in later burials, it raises the possibility as to whether their presence indicates a closer connection with the earlier phases of the Napatan period.


The sole nature-related motif was of a flower. Included in the “other” category are four pierced nuggets, three of gold and one amazonite. There was no indication that they were inscribed or decorated. Furthermore, three amulet cases were found, but, in contrast to the one from the 25th Dynasty-Early Napatan tomb, they were all hollow.

 \begin{figure}[H]
	\centering
	\begin{minipage}{0.49\textwidth}
		\centering
		\includegraphics[width=\textwidth, height=4cm, keepaspectratio]{25-MN-23.831}
	\end{minipage}
	\hspace{-25mm}%
	\begin{minipage}{0.49\textwidth}
		\centering
		\includegraphics[width=\textwidth, height=4cm, keepaspectratio, clip, trim={10cm 10cm 10cm 15cm}]{25-MN-23.349}
	\end{minipage}
	
	\begin{minipage}{0.49\textwidth}
		\centering
		\includegraphics[width=\textwidth, height=4cm, keepaspectratio, clip, trim={6cm 5cm 5cm 6cm}]{25-MN-47.3916}
	\end{minipage}	
	\hspace{-30mm}%
	\begin{minipage}{0.49\textwidth}
		\centering
		\includegraphics[width=\textwidth, height=4cm, keepaspectratio, clip, trim={5cm 5cm 5cm 5cm}]{25-MN-47.3903}
	\end{minipage}
	\caption{Original necklace of 31 Bes and one sow amulet; amulets of \textit{udjat}, monkey and sow (MFA 23.831, 47.3916, 23.349, 47.3903)}
	\label{fig:25-MN-pics}
\end{figure}


\paragraph{Amulet materials}


\begin{figure}[H]
	\centering
	\includegraphics[width=1\textwidth]{25-MN-materials}
	\caption{Distribution of amulet materials in 25th Dynasty-Middle Napatan non-elite tombs.}
	\label{fig:25-MN-materials}
\end{figure}


The amulets were made of only five different materials: faience, silver, gold, chalcedony and amazonite (Fig. 134). Faience continued to be the most common material, comprising 91\% of all amulets. Silver, chalcedony and amazonite, materials likely sourced from the Egyptian territory of Lower Nubia, make up 6\% of amulets, with silver being more common. All of the silver amulets were either amulet cases or \textit{udjat}. The sole chalcedony example was of the monkey, while the amazonite amulet was the pierced nugget. Gold amulets comprise only 3\% of the total, and include an \textit{udjat} and the pierced nuggets. This makes the number of gold amulets in this phase significantly lower than the 25th Dynasty-to-Early Napatan assemblages.


\subsection{Early Napatan and Early Napatan-Middle Napatan}


For the purpose of this analysis, Early and Middle Napatan non-elite graves are considered as a single group. This consolidation is necessary given the small Early-Middle Napatan sample, which comprises 15 tombs and only one amulet.


\subsubsection{Tomb structure}


Early Napatan non-elite tombs were more concentrated in Meroe South, with a total of 43 tombs, whereas only one in Meroe West was assigned this dating (Figure \ref{fig:EN-EN-MN-tombs}). All substructures were pits and superstructures were classified as either non-existent or unidentified. The same happens with Early-Middle Napatan tombs, although they are less numerous, with 13 tombs in Meroe South and two in Meroe West. The only tomb with an additional architectural feature is S 44 from Early Napatan, which had a unidentified superstructure made of masonry and a slope as a form of entrance into the burial pit. Parallels to this feature are eight Meroe South elite graves from the 25th Dynasty-to-Middle Napatan and nine non-elite tombs from the 25th Dynasty-Early Napatan phase at Sanam.

\begin{figure}[H]
	\centering
	\includegraphics[width=1\textwidth]{EN-EN-MN-tombs}
	\caption{Distribution of super- and substructures in Early Napatan and Early-Middle Napatan non-elite tombs.}
	\label{fig:EN-EN-MN-tombs}
\end{figure}

\subsubsection{Object types}


The tombs at Meroe West from both phases included very few grave goods: Early Napatan tomb W 616 contained only two ceramic vessels, while W 584 from Early-Middle Napatan included two beads and three ceramic vessels (Figure \ref{fig:EN-EN-MN-objs}). Only five Meroe South tombs from Early-Middle Napatan contained objects, among them three bead nets, a decorative element, a non-ceramic vessel, 71 loose beads and one ceramic vessel. Considering the scarcity of grave goods, the presence of bead nets is notable as the sole burial-specific object, which was included already on previous non-elite contexts. Traces of a coffin were found in ten graves: two from Early Napatan and eight from Early-Middle Napatan. Similar to shabtis, all of these tombs were at Meroe South.


Early Napatan tombs from Meroe South are the most diverse in terms of grave goods, where only 20\% of graves did not contain any objects. For the most part, the assemblage is very similar to those that preceded it. However, with one difference: the number of shabtis found in two burials. S 107 contained only two shabtis, which were noted by Reisner as possibly intrusive \parencites[p.~426]{dunhamWestSouthCemeteries1963}. Conversely, S 132 had 117 shabtis, with some of them found dispersed in other graves by the looters. 110 of these figures were inscribed with the name Sesensa and were similar in style as those from Tanwetamani \parencites[p.~362]{dunhamWestSouthCemeteries1963}. If, indeed, S 132 was the source of all the shabtis, this particular tomb is an exceptional case. There is no other indication in the architecture and other artifacts indicating that it was the burial of a high status individual. This evidence leads to the possibility that either the deceased or their relatives/community were able to acquire a significant quantity of shabtis, similar in style to royal ones, and have them inscribed―if, indeed, the name in the shabtis is of the deceased. Furthermore, all non-elite shabtis were found in Meroe South burials (see sections above), suggesting that the region provided a level of access to royal-style objects that did not exist in other non-elite contexts. However, as most shabtis from Meroe South tombs were found scattered and the site also contained Meroitic graves, we can not discard the likelihood that these shabtis could belong to these later burials.

\begin{figure}[H]
	\centering
	\begin{minipage}{1\textwidth}
		\centering
		\includegraphics[width=\textwidth, keepaspectratio]{EN-objs}
	\end{minipage}
	
	\begin{minipage}{1\textwidth}
		\centering
		\includegraphics[width=\textwidth, height=6cm, keepaspectratio]{EN-MN-objs}
	\end{minipage}
	\caption{Distribution of object types in Early Napatan and Early-Middle Napatan non-elite tombs.}
	\label{fig:EN-EN-MN-objs}
\end{figure}


\subsubsection{Amulets}


\paragraph{Amulet distribution}


{\small
	\begin{xltabular}{\textwidth}{X X X X X}
		\caption{Distribution of amulets in Early Napatan and Early-Middle Napatan non-elite tombs by site, including total counts of amulets and tombs, tombs containing amulets and amulet frequency (\%).}
		\label{tab:EN-EN-MN-distro} \\
		\toprule
		\textbf{Site} & \textbf{Total number of tombs} & \textbf{Total number of amulets} & \textbf{N. of tombs with amulets} & \textbf{\% of tombs with amulets} \\
		\midrule
		\endfirsthead
		\multicolumn{5}{c}{\textit{Continued from previous page}} \\
		\toprule
		\textbf{Site} & \textbf{Total number of tombs} & \textbf{Total number of amulets} & \textbf{N. of tombs with amulets} & \textbf{\% of tombs with amulets} \\
		\midrule
		\endhead
		\bottomrule
		\endlastfoot
		Meroe South (EN) & 43 & 25 & 7 & 28\% \\
		\cmidrule(r){1-5}
		Meroe West (EN--MN) & 2 & 1 & 1 & 100\% \\
		\midrule
		\textbf{Total} & \textbf{45} & \textbf{26} & \textbf{8} & -- \\
	\end{xltabular}
}


Amulets from Early Napatan were found in seven Meroe South tombs, 28\% of the total graves. 25 amulets were distributed across seven burials (Table \ref{tab:EN-EN-MN-distro}). The sole Early-Middle Napatan amulet was from a tomb in Meroe West. This is compatible with the trend attested in royal and elite contexts, where the number of amulets decrease in later phases of the Napatan period.


\paragraph{Amulet types}

\begin{figure}[H]
	\centering
	\includegraphics[width=1\linewidth, height=6cm, keepaspectratio]{EN-EN-MN-types}
	\caption{Distribution of amulet types in Early Napatan and Early-Middle Napatan non-elite tombs.}
	\label{fig:EN-EN-MN-types}
\end{figure}


Early Napatan amulets are divided into five types, with symbols being the most common, making up 80\% of the total amulets. This distribution is more similar to royal assemblages than previous non-elite ones. Two amulets were of miscellaneous motifs, although both were unable to be identified, and there was only one instance each of animal, deity and object motifs. The only Early-Middle Napatan amulet was of a symbol (Figure \ref{fig:EN-EN-MN-types}).


\paragraph{Amulet motifs}


\begin{figure}[H]
	\centering
	\includegraphics[width=1\textwidth, height=6cm, keepaspectratio]{EN-EN-MN-motifs}
	\caption{Distribution of amulet motifs in Early Napatan and Early-Middle Napatan non-elite tombs.}
	\label{fig:EN-EN-MN-motifs}
\end{figure}


\begin{figure}[H]
	\centering
	\begin{minipage}{0.25\textwidth}
		\centering
		\includegraphics[width=\textwidth, height=5cm, keepaspectratio]{EN-EN-MN-21.12192}
	\end{minipage}
	\hspace{-5mm}%
	\begin{minipage}{0.25\textwidth}
		\centering
		\includegraphics[width=\textwidth, height=5cm, keepaspectratio]{EN-EN-MN-21.12252.1}
	\end{minipage}
	\hspace{-5mm}%
	\begin{minipage}{0.15\textwidth}
		\centering
		\includegraphics[width=\textwidth, height=3.5cm, keepaspectratio, clip, trim={8cm 8cm 8cm 7cm}]{EN-EN-MN-21.12212}
	\end{minipage}
	\caption{Early Napatan amulets: two \textit{udjat} (MFA 21.12192, 21.12252.1), and two-headed monkey (MFA 21.12212).}
	\label{fig:EN-EN-MN-pics}
\end{figure}


Among the Early Napatan amulets, the \textit{udjat} was the most popular (76\%) (Figure \ref{fig:EN-EN-MN-motifs}). Different designs were noted, such as carved and openwork (Figure \ref{fig:EN-EN-MN-pics}). The only oyher symbol amulet was of a quadruple \textit{udjat}, a motif very popular in 25th Dynasty-Early Napatan assemblages. The sole Early-Middle Napatan amulet was of the \textit{udjat}.


The only deity amulet was, unfortunately, unidentified. A gold axe head amulet was found as the sole object in tomb S 101A. This tomb was an intact burial pit, but included a coffin and traces of mummification on the body. Since this was an undisturbed burial, the evidence indicates an interesting interplay of social status markers. Although the tomb architecture was very simple, the deceased or family members/community were able to have access to a coffin, a process to preserve the body, and, albeit there was only this single amulet as a grave good, it was made of gold and of a motif not found elsewhere, meaning it was not produced en masse.


Another interesting amulet is of a monkey/baboon, likely representing the god Thoth (Figure \ref{fig:EN-EN-MN-pics}). Even though it was not noted neither on the report nor in the MFA's description, it seems like the baboon has two heads. If true, this amulet would be the second one with a double head, the first being the unidentified deity with two hawk heads from a royal tomb from the 25th Dynasty.


\paragraph{Amulet materials}


\begin{figure}[H]
	\centering
	\includegraphics[width=1\textwidth, height=6cm, keepaspectratio]{EN-EN-MN-materials}
	\caption{Distribution of amulet materials in Early Napatan and Early-Middle Napatan non-elite tombs.}
	\label{fig:EN-EN-MN-materials}
\end{figure}


The only Early-Middle Napatan \textit{udjat} amulet was made of faience, which as also the most common material for Early Napatan amulets (48\%) (Figure \ref{fig:EN-EN-MN-materials}). Similar to the 25th Dynasty-Middle Napatan phase, silver amulets are the second most popular, but in this phase in higher quantity (40\%). There was only one gold amulet as well, the one depicting the axe head. Carnelian and amazonite were the composition of two amulets each. Both stones were not sourced from the Nubian region, carnelian most likely came from a mine near Aswan and, amazonite, from either Lower Nubia or Ethiopia. Although imported materials appear to be the majority in Early Napatan amulets, it is important to consider that silver could also be obtained by recycling older artifacts.


\section{Regional variations}


Given the geographical spread of the non-elite sites across Nubia, this study also employs a regional framework to analyze variation in amulet assemblages. The sites are categorized into four main centers, from north to south: Lower Nubia (Mirgissa, Qustul and Debeira 176), north Upper Nubia (Missiminia), Fourth Cataract (Sanam), Meroe region (Meroe West and South). This approach aims to understand if a greater distance with the northern territory influences the use of certain amulets.


Distribution of amulets and tombs per region is irregular. Although Lower and North Upper Nubia contain less amulets than tombs, the Fourth Cataract and Meroe region show the opposite (Figure \ref{fig:regional-distro}). Furthermore, the Fourth Cataract, despite encompassing only one site, is the region that yielded the highest total number of amulets. However, amulets were found in small minority of tombs across all regions, with none exceeding 23\%. The Fourth Cataract and Meroe regions had the highest proportion (23\%), followed by North Upper Nubia (13\%) and Lower Nubia (9\%).

\begin{figure}[H]
	\centering
	\includegraphics[width=1\textwidth]{regional-distro}
	\caption{Distribution of total number of non-lite tombs and amulets by region.}
	\label{fig:regional-distro}
\end{figure}


\begin{figure}[H]
	\centering
	\includegraphics[width=1\textwidth]{regional-types}
	\caption{Distribution of non-elite amulet types by region.}
	\label{fig:regional-types}
\end{figure}


Concerning the distribution of amulet types, Lower Nubia had the lowest concentration of all types, with deities and symbols being slightly more common (Figure \ref{fig:regional-types}). In contrast, the spread of deity amulets at North Upper Nubia was higher, but the region lacked any amulets of objects, humans, nature and miscellaneous types. The Fourth Cataract region exhibits a significant increase in the number of amulets depicting symbols, deities and animals, while human, object, and nature motifs did not increase in the same pace, being slightly higher than the prior region. At the Meroe region, although symbol and human amulets increase moderately, deity and animal motifs decrease considerably. Nature amulets also drop slightly, while object and miscellaneous motifs remain virtually the same.


\begin{figure}[H]
	\centering
	\includegraphics[width=1\textwidth]{regional-udjat}
	\caption{Distribution of non-elite \textit{udjat} and quadruple \textit{udjat} amulets by region.}
	\label{fig:regional-udjat}
\end{figure}


The \textit{udjat}, with 1878 total instances, is the most numerous motif across all non-elite contexts. A regional analysis of its distribution reveals that the notable increase in symbol amulets within the Meroe region was driven almost entirely by this single motif, which constitutes nearly 90\% of all symbol amulets there (Figure \ref{fig:regional-udjat}). Quadruple \textit{udjat} were not as recurrent in most sites, with the exception of North Upper Nubia where this motif was almost two times more frequent than the normal \textit{udjat}. Furthermore, at Lower Nubia, the different quantity of both amulets was not particularly substantial like it was at the Meroe region and Fourth Cataract.


Upon analyzing exclusive amulet motifs to each region, the Fourth Cataract emerges once again as the most diverse (Figure \ref{fig:regional-motifs}). This location yielded 34 motifs not found in any other non-elite site. Among those, the majority represent deities and symbols (74\%), although all types of amulets are represented. The Meroe region yielded 24 unique motifs, with most representing deities, animals and humans (71\%). Only three symbol amulets were exclusive to North Upper Nubia, while all motifs found in Lower Nubia were found elsewhere.


The regional distribution of local and imported materials provides another key metric (Figure \ref{fig:regional-materials}). While locally sourced materials were predominant everywhere, variations exist. The Meroe region, for instance, had the highest proportion of local materials and the lowest of imports. This supports the hypothesis that materials imported from Egypt became less accessible further south of the Nubian territory. Interestingly, despite Sanam's role as a trade hub, it yielded fewer imported materials than the more northerly region of North Upper Nubia, underscoring the importance of proximity to Egypt. However, the relatively low distribution of imported materials at Lower Nubia suggests that proximity to Egypt was not the sole factor regulating access, and that other social or economic dynamics were likely at play.


\begin{figure}[H]
	\centering
	\includegraphics[width=1\textwidth, height=12cm, keepaspectratio]{regional-motifs}
	\caption{Exclusive non-elite amulet motifs by region.}
	\label{fig:regional-motifs}
\end{figure}


\begin{figure}[H]
	\centering
	\includegraphics[width=1\textwidth]{regional-materials}
	\caption{Distribution of imported and local amulet materials in non-elite tombs by region (in \%).}
	\label{fig:regional-materials}
\end{figure}


\section{Conclusion}


This chapter had presented and analyzed data from elite and non-elite cemeteries of the Napatan period across temporal and regional dimensions. The analysis of tomb structure revealed burials belonging to individuals of a higher status. The grave good assemblages largely conform to the established Napatan patterns, although the 25th Dynasty to Early Napatan phases introduced new, albeit rare, object types. Additionally, objects manufactured specifically for the burial were less common on non-elite contexts than in royal or elite tombs. Shabtis, typically considered a royal-exclusive object during the Napatan period, were found in tombs of the 25th Dynasty and Early Napata phases. However, most were intrusive fragments that could have originated in an elite context. The only shabtis that Reisner/Dunham attributed to a probable origin to are the 110 inscribed fragments from the Early Napatan tomb S 132. If this attribution is correct, this tomb constitutes a significant exception, indicating that its occupant or community had the means to acquire personalized shabtis despite the tomb's non-elite character.


The amulet analysis revealed parallel trends	between all social groups: both non-elite and elite assemblages show a peak in quantity and diversity during the 25th Dynasty to Early Napatan period, succeeding the low number at pre-25th Dynasty and followed by a decline in later phases. However, the non-elite sphere displays distinct characteristics, including a greater variety of animal motifs and unique adaptations of Egyptian-style designs. Plaques and amulets combining multiple motifs were also more prevalent in non-elite contexts. While symbol amulets dominated non-elite graves in the 25th Dynasty to Early Napatan phases, deities, animals and humans were more popular in the preceding and subsequent periods. Material use also shifted temporarily, with an increase in imported materials during the 25th Dynasty to Early Napatan period. However, local materials remained dominant overall, driven mostly by faience, the most popular material composition across all phases.


Regionally, the Fourth Cataract contained the highest total number of amulets and graves, and yielded the most unique designs, particularly of deities and symbols. In contrast, the Meroe region had largest concentration of symbol and human motifs, especially the \textit{udjat} and depictions of boys. Both regions, however, contained the same proportion of tombs with amulets. The use of imported materials was more pronounced in the northern regions, consistent with their primary source in the Egyptian Eastern Desert.
In the next chapter, this temporal and spatial analysis will be integrated with the distinct social contexts in order to build a multifaceted interpretation of Napatan amulets.